\chapter{方程式の品質保証のための基本定理}\label{chap:basic_theorem}

%
%----------------------------------------------------------------------------------------------------------------------------------------------
%----------------------------------------------------------------------------------------------------------------------------------------------
%----------------------------------------------------------------------------------------------------------------------------------------------
\section{Banach空間}\label{sec:Banach_space}
Banach空間を知らないという場合，実数全体の集合${\mathbb R}$や複素数全体の集合${\mathbb C}$の直積空間${\mathbb R}^{N}$や${\mathbb C}^N$，はたまた${\mathbb R}$あるいは${\mathbb C}$を表す記号${\mathbb K}$と読み替えても，ほとんどの章は差し支えありません．
では，なぜBanach空間で紹介するのか?という疑問が生まれると思います．
偏微分方程式の解は一般的に無限次元となり，${\mathbb R}^{N}$では表せません．
このような問題を扱う際には無限次元まで扱えるBanach空間が必要になります．
${\mathbb R}^{N}$もBanach空間であるため，連立一次方程式の解の品質保証を行いたい場合は${\mathbb R}^{N}$で十分です．

%----------------------------------------------------------------------------------------------------------------------------------------------
\subsection{Banach空間の定義と性質}\label{sec:define_Banach}
みなさんは有理数${\mathbb Q}$と実数${\mathbb R}$の違いを明確に答えられますか?
誤解を恐れずにいうと「有理数に抜け目のないように穴を埋めたものが実数」といえます．
すなわち，実数の上を歩いたとしても落とし穴にはまることがない集合，完備な集合といえます．
実数の完備性は色々な出発点と定理があるが，その中に「Cauchy列ならば収束列」であることを示す出発点あるいは定理があったことを思い出しておくと，すんなりとこの節を読むことができると思います．


この節では，実数に限らず一般的な線形空間でも，抜け目がない集合を定義することが目的です．
この抜け目のない線形空間はBanach空間と呼ばれます．
非常に抽象的な定義で難しいと感じるかもしれませんが，実数の完備性を思い出すと，実は一般化しただけの自然にも思える定義に感じれると思います．

では，線形空間がないと話が進みませんので，線形空間の公理から始めましょう．

\begin{axiom}[線形空間の公理]\label{axiom:linear_space}
空でない集合$X$が係数体${\mathbb K}$上の線形空間であるとは，
任意の$u, v \in X$とスカラー$\alpha \in {\mathbb K}$に対して，加法$u + v \in X$とスカラー乗法$\alpha u \in X$が定義されていて，
任意の$u,v,w \in X$とスカラー$\alpha, \beta \in {\mathbb K}$に対して次の(i)-(viii)が成り立つことである．\\
(i) $(u + v) + w = u + (v + w)$ \\
(ii) $u + v = v + u$ \\
(iii) $u + 0 = u$となる$0 \in X$が一意に存在 \\
(iv) $u + (-u) = 0$となる$-u \in X$が一意に存在 \\
(v) $\alpha (u + v) = \alpha u + \alpha v$ \\
(vi) $( \alpha + \beta ) u = \alpha u + \beta u$ \\
(vii) $( \alpha \beta ) u = \alpha (\beta u)$ \\
(viii) $1 u = u, ~ 1 \in {\mathbb K}$
\end{axiom}

再び，実数の完備性を思い出してみましょう．
実数の場合は，真っ先に$\varepsilon - N$論法を使って数列の収束を勉強したと思います．
この際に，自然と絶対値記号を使っていたと思います．
しかし，一般的な線形空間には絶対値記号はありません．
例えば，${\mathbb C}^2$の元
\begin{align*}
	\left( \begin{array}{c}
		2 + 2 i \\
		-1 + 3 i
	\end{array} \right)
\end{align*}
に対する「絶対値に相当するもの」は何でしょうか?
ここで重要なのは，あくまで一般的な線形空間$X$の点列に対して$\varepsilon - N$論法を使って収束を定義したいことに着目すると「絶対値に相当するもの」を${\mathbb C}^2$の元につけた結果が${\mathbb C}^2$では$\varepsilon - N$論法の$\varepsilon$がよくわからない状態になってしまいますよね．
そのために，どんな線形空間$X$の元でも「絶対値に相当するもの」をつけた結果は実数${\mathbb R}$になるように定義します．
もちろん，「絶対値に相当するもの」であるため，絶対値の性質も含めて定義する必要もあります．
これらの点に留意しながら次の「絶対値に相当するもの」の定義，すなわちノルムの定義を見てみましょう．

\begin{defi}[ノルムとノルム空間の定義]\label{def:App:norm}
$X$を係数体$\mathbb{K}$上の線形空間とする．
$X$で定義された関数$\|\cdot\|: X \rightarrow {\mathbb R} $が$X$のノルムであるとは\\
(i) $\| u \| \ge 0, ~ u\in X$ \\
(ii) $\| u \|=0 \Leftrightarrow u=0$\\
(iii) $\| \alpha u \| = |\alpha| \| u \|, ~ (\alpha \in {\mathbb K},~ u \in X)$\\
(iv) $\| u+v \| \le \|u\| + \|v\|\\$
が成立することである．
さらに$X$に一つのノルムが指定されているとき，$X$はノルム空間という．
\end{defi}

これで，線形空間にノルムという絶対値に相当する道具が追加されたノルム空間が定義できました．
まだ，ノルム空間が$X$しか出てきていないので$\| \cdot \|$でも問題ありませんが，ノルム空間が多数出てきた場合，$\| \cdot \|_{X}$のように右下にどのノルム空間のノルムなのか明記します．
続いてノルムを使って，点列の収束を定義してみましょう．


\begin{defi}[ノルム空間の収束と極限]\label{def:App:Con}
$X$をノルム空間とする．
$X$の点列$(u_n) \subset X$は
\begin{align*}
	\forall \varepsilon > 0, ~ \exists N \in {\mathbb N}, ~ \forall n \ge N ~ \mbox{に対して} ~ \| u_n - u \| < \varepsilon
\end{align*}
のとき，点$u \in X$に収束するといい，
\begin{align*}
\| u_n - u \| \rightarrow 0, ~ (n \rightarrow \infty)
\end{align*}
と表す．
このとき，$u$を$u_n$の極限といい
\begin{align*}
u_n \rightarrow u, ~ (n\rightarrow \infty)
\end{align*}
と表す．
\end{defi}

\begin{Thm}[極限の一意性]\label{thm:limit_unique}
	$X$の点列$(u_n) \subset X$が収束するならば，その極限は一意である．
\end{Thm}

\begin{Proof}
極限を$v$と$w$の2つ持つと仮定して矛盾を示す．
つまり，$v$と$w$が極限であることから，
\begin{align*}
	\forall \varepsilon > 0, ~ \exists N_1 \in {\mathbb N}, ~ \forall n \ge N_1 ~ \mbox{に対して} ~ \| u_n - v \| < \varepsilon
\end{align*}
と
\begin{align*}
	\forall \varepsilon > 0, ~ \exists N_2 \in {\mathbb N}, ~ \forall n \ge N_2 ~ \mbox{に対して} ~ \| u_n - w \| < \varepsilon
\end{align*}
が成立する．
そのうえで，$\epsilon$は任意であることから，
\begin{align*}
	\epsilon = \frac{ \| v - w \|  }{2}
\end{align*}
としたとき，上の論理式を満たす$N_1$と$N_2$がそれぞれ存在する．
すなわち，
\begin{align*}
	\exists N_1 \in {\mathbb N} ~ \forall n \ge N_1 ~ \mbox{に対して} ~ \| u_n - v \| < \frac{ \| v - w \|  }{2}
\end{align*}
と
\begin{align*}
	\exists N_2 \in {\mathbb N} ~ \forall n \ge N_2 ~ \mbox{に対して} ~ \| u_n - w \| < \frac{ \| v - w \|  }{2}
\end{align*}
となる．
そのうえで，$N_{\max} = \max (N_1, N_2)$とすると，
\begin{align*}
	\forall n \ge N_{\max} ~ \mbox{に対して} ~ \| u_n - v \| < \frac{ \| v - w \|  }{2} ~~ \mbox{および} ~~ \| u_n - w \| < \frac{ \| v - w \|  }{2}
\end{align*}
となる．
よって，任意の$n \ge N_{\max}$に対して，
\begin{align*}
	\| v - w \| \le \| v - u_n \| + \| u_n - w \| < \frac{ \| v - w \|  }{2} + \frac{ \| v - w \|  }{2} = \| v - w \|
\end{align*}
となり，$\| v - w \|  <  \| v - w \|$から，矛盾する．

	\qed
\end{Proof}


同様にノルム空間上のCauchy列もノルムを用いて次のように定義できます．

\begin{defi}[Cauchy列]\label{def:App:Cau}
$X$をノルム空間とする．
そのとき$X$の点列$(u_n)$がCauchy列であるとは
\begin{align*}
u_n - u_m \rightarrow 0, ~ (n,m \rightarrow \infty)
\end{align*}
が成立することである．
即ち
\begin{align*}
\| u_n - u_m \| \ \rightarrow 0, ~  (n,m \rightarrow \infty)
\end{align*}
が成立することである．
\end{defi}

ちなみに$\varepsilon - N$論法を用いてちゃんと書くと，$X$の点列$(u_n)$が
\begin{align*}
  \forall \varepsilon > 0, ~ \exists N \in {\mathbb N}, ~ \forall n, m \ge N ~ \mbox{に対して} ~ \| u_n - u_m \| < \varepsilon
\end{align*}
を満たすときに点列$(u_n)$がCauchy列であるといいます．

点列の収束とCauchy列は一見すると似ているように見えますが，明確に意図の違いがあります．
点列の収束の定義は，$X$の点列$(u_n)$の極限$u$が$X$に属していることが既にわかっています．
しかし，抽象的なノルム空間$X$の点列$(u_n)$の極限が$X$に属しているか，それとも属していないかなんてわかりません．
そのために，極限を使わずに収束してそうな点列としてCauchy列が定義されたと思っておいてください．
もちろん，ノルム空間$X$の収束する点列$(u_n)$ならば，収束しそうな点列であるCauchy列になることは簡単にいえます．
実際に，収束する点列$(u_n) \subset X$は極限$u \in X$を持ちます．その上，
\begin{align*}
&\| u_n - u_m \| = \| u_n - u + u - u_m \| \\
&\le \| u_n - u \| + \| u - u_m \| \rightarrow 0, ~ (n, m \rightarrow \infty).
\end{align*}
から，収束する点列$(u_n) \subset X$ならば，Cauchy列になります．

%もちろんノルム空間$X$上のCauchy列だが，極限を$X$に持たない例もあります．
%例えば，有理数全体の集合${\mathbb Q}$は体になるので，係数体${\mathbb K}= {\mathbb Q}$, $X = Q$として，絶対値を導入することでノルム空間になります．
%その上，連分数を用いて円周率$\pi$に収束するような${\mathbb Q}$上の点列を考えてみてれば一目瞭然ですよね．

Cauchy列が収束しそうな点列であることを踏まえたうえで，次の完備に関する定義を見てみましょう．

\begin{defi}[完備]\label{def:App:Comp}
$X$をノルム空間とする．
$X$が完備であるとは，任意のCauchy列$(u_n)$が$X$の中で極限を持つことである．
すなわち，任意のCauchy列$(u_n) \subset X$が
\begin{align*}
\| u_n - u \| \rightarrow 0, ~ (n \rightarrow \infty)
\end{align*}
となる極限$u$を$X$内に持つことである．
\end{defi}

ノルム空間$X$が完備であるとき，どんなCauchy列$(u_n)$でも，ちゃんと極限$u$を$X$内に持つことを保証してくれています．
この節の目的は「抜け目のない線形空間」を定義することでしたね．
ノルム空間が完備であれば抜け目のない線形空間の出来上がりです．
すなわち，Banach空間の定義は次のようになります．

\begin{defi}[Banach空間]\label{def:App:Banach}
ノルム空間$X$がBanach空間であるとは，$X$が完備であることである．
\end{defi}

この節の最後に，後々必要になるCauchy列の性質を示しておきます．
まず，Cauchy列の性質を示すために必要になるノルムの性質として逆三角不等式を示します:

\begin{Thm}[逆三角不等式]\label{thm:inv_triangle_inq}
$X$をノルム空間とする．
任意の$u, v \in X$について次の不等式を満たす:
\begin{align*}
| \| u \| - \| v \| | \le \| u - v \|
\end{align*}
\end{Thm}

\begin{Proof}
任意の$u, v \in X$について
\begin{align*}
 \| u \| &= \| u - v + v \| \le \| u - v \| + \| v \| \\
 \| v \| &= \| v - u + u \| \le \| v - u \| + \| u \| = \| u - v \| + \| u \| 
\end{align*}
となる．
よって
\begin{align*}
 \| u \| - \| v \| &\le \| u - v \| \\
 \| v \| - \| u \|  &\le  \| u - v \| 
\end{align*}
となるため，
\begin{align*}
| \| u \| - \| v \| | \le \| u - v \|
\end{align*}
を持つ．

\qed
\end{Proof}

続いて，Cauchy列とは別の点列の性質を定義します:

\begin{defi}[有界列]\label{def:bounded_seq}
$X$をノルム空間とする．
そのとき$X$の点列$(u_n)$が有界列とは任意の$n \in {\mathbb N}$に対して
\begin{align*}
\| u_n \| \le M
\end{align*}
となる定数$M > 0$が存在することである．
\end{defi}

有界列のポイントは点列$(u_n)$に対して，1つの$M$が存在することです．
すなわち，有界列なら，どんな$u_n$でも共通の$M$を用いて$\| u_n \| \le M$という不等式を満たします．

\begin{Thm}[Cauchy列ならば有界列]\label{def:if_cauchy_then_bounded}
$X$をノルム空間とする．
そのとき$X$の点列$(u_n)$がCauchy列ならば有界列でもある．
\end{Thm}

\begin{Proof}
$X$の点列$(u_n)$がCauchy列であるため，$\varepsilon - N$論法を用いた表記で
\begin{align*}
  \forall \varepsilon > 0, ~ \exists N \in {\mathbb N}, ~ \forall n, m \ge N ~ \mbox{に対して} ~ \| u_n - u_m \| < \varepsilon
\end{align*}
を満たす．
$\varepsilon = 1$としても，それに対応した$N$が存在し，任意の$n \ge N$に対して
\begin{align*}
\| u_n - u_N \| < 1
\end{align*}
を満たす．(ここで$m$を$N$に固定したことに注意．)

任意の$n \ge N$に対して$\| u_n \|$が$\| u_N \|$で評価できることを示す．
逆三角不等式である定理\ref{thm:inv_triangle_inq}を用いると
\begin{align*}
| \| u_n \| - \| u_N \| | \le  \| u_n - u_N \| < 1
\end{align*}
となる．
絶対値の性質より$| \| u_n \| - \| u_N \| | < 1$は
\begin{align*}
\| u_N \| - 1 < \| u_n \| < \| u_N \| + 1
\end{align*}
となる．
よって
\begin{align*}
M = \max \{ \| u_1 \|, ~ \| u_2 \|, ~ \cdots, ~ \| u_{N -1} \|, ~ \| u_N \| + 1 \}
\end{align*}
とすると，任意の$n \in {\mathbb N}$について
\begin{align*}
\| u_n \| \le M
\end{align*}
が成り立つため，点列$(u_n)$は有界列である．


\qed
\end{Proof}



%----------------------------------------------------------------------------------------------------------------------------------------------
\subsection{ノルム空間の位相構造とBanach空間の閉部分空間}\label{sec:subset_Banach}
まず，線形代数の復習である線形部分空間を定義しましょう．

\begin{defi}[線形部分空間]\label{def:linearsubspace}
	線形空間$X$の空ではない集合$M$が任意の元$u,v \in M$と任意の係数体$\alpha \in {\mathbb K}$に対して
	\begin{align*}
	u + v \in M \\
	\alpha u \in M
	\end{align*}
	を満たすとき，$M$は線形空間$X$の線形部分空間と呼ぶ．
\end{defi}

定義\ref{def:linearsubspace}の線形部分空間は線形代数で習っている通り，$M$も$X$の演算のもとに線形空間となります．
では，Banach空間の線形部分空間は，必ずBanach空間になるでしょうか?
答えは残念ながらNoです．
Banach空間の線形部分空間はBanach空間にはならず，完備性がないノルム空間になってしまう場合もあります．
では，Banach空間の線形部分空間がBanach空間になるための条件は何でしょうか?
本節では，この条件を紹介することが目標です．
その際に定義するノルム空間の位相構造も，距離空間における集合と位相の復習となり，本節以外にも出てくるので覚えておきましょう．
まず，ノルム空間における開球，開集合，閉集合を定義します．

\begin{defi}[ノルム空間の開球]\label{def:App:OpenBall}
	$X$をノルム空間とする．
	$x \in X$とし，$r > 0$を正実数とする．
	そのとき，集合
	\begin{align*}
		B_X(x, r) := \{ y \in X ~ | ~ \| x - y \|_{X} < r \}
	\end{align*}
	を中心$x$，半径$r$の開球という．
	$X$が明らかな場合は$B_X(x, r)$を省略して$B(x, r)$と表記する．
\end{defi}

\begin{defi}[ノルム空間の開集合]\label{def:App:OpenSet}
	$X$をノルム空間し，$M$を$X$の部分集合とする．
	任意の$x \in M$に対して，$B_X(x, r) \subset M$となる$r > 0$が存在する場合，$M$が開集合であるという．
\end{defi}

\begin{defi}[ノルム空間の閉集合]\label{def:App:CloseSet}
	$X$をノルム空間とし，$M$を$X$の部分集合とする．
	$M$が閉集合であるとは，$M$の任意の点列$(u_n)$の極限$u \in X$が$M$にも属することである．
	すなわち，点列$(u_n) \subset M$について
	\begin{align*}
		u_n \rightarrow u, ~(n \rightarrow \infty) \Rightarrow u \in M 
	\end{align*}
	であるとき，$M$は閉集合であるという．
\end{defi}

集合と位相の講義で習った一般的な位相空間$X$の部分集合$M$の閉集合の定義は補集合
\begin{align*}
	M^{c} := \{ x \in X ~ | ~ x \not\in M \}
\end{align*}
としたとき，$M^{c}$が開集合になることでしたね．
ノルム空間の場合は，「$M^{c}$が開集合」と「定義\ref{def:App:CloseSet}」は同値になります．
そのため，本書では，良く使う定義\ref{def:App:CloseSet}を閉集合の定義としました．
もちろん，「$M^{c}$が開集合」を閉集合の定義として，定理として定義\ref{def:App:CloseSet}と同値であることを導く方針もあります．


それでは本節の目標である「Banach空間の線形部分空間がBanach空間になるための条件」を考えていきます．
先に答えを言ってしまうと，「Banach空間の線形部分空間がBanach空間になるための条件」は次に定義する「閉部分空間」と同値になります．
では閉部分空間について定義しましょう．

\begin{defi}[閉部分空間]\label{def:App:Closed_liner_subspace}
	$X$をノルム空間とし，$M$を$X$の線形部分空間が閉集合であるとき，$M$を閉部分空間であるという．
\end{defi}

端的に言えば，「線形部分空間」かつ「閉集合」のとき，「閉部分空間」と呼ばれます．
さて，いよいよ，「Banach空間の線形部分空間がBanach空間になるための条件」を示しましょう．

\begin{Thm}[線形部分空間とBanach空間の関係]\label{thm:Banach_space_closed_subspace}
$X$をBanach空間とし，$M$を$X$の線形部分空間とする．
そのとき，
\begin{align*}
	M\mbox{が}X\mbox{の閉部分空間} ~ \Leftrightarrow ~ M\mbox{が}X\mbox{のノルムでBanach空間}
\end{align*}
\end{Thm}

\begin{Proof}
「$M\mbox{が}X\mbox{の閉部分空間} ~ \Rightarrow ~ M\mbox{が}X\mbox{のノルムでBanach空間}$の証明」

点列$(u_n)$を$M$の任意のCauchy列とする．
$M$はBanach空間$X$の部分集合であることから，Banach空間$X$の完備性(定義\ref{def:App:Comp}と定義\ref{def:App:Banach})より，$M$の任意のCauchy列$(u_n)$の極限$u$が$X$内に存在する．
さらに，$M$が$X$の閉部分空間であることから，閉集合の定義\ref{def:App:CloseSet}より任意の点列の極限は$M$に属するため，Cauchy列$(u_n)$の極限$u$は$M$にも属する．
よって，ノルム空間$M$の任意のCauchy列$(u_n)$の極限$u$が$M$に属するため，$M$はBanach空間となる．\\

\noindent 「$M\mbox{が}X\mbox{の閉部分空間} ~ \Leftarrow ~ M\mbox{が}X\mbox{のノルムでBanach空間}$の証明」

$M$の任意の点列を$(u_n)$とすると，$u_n \rightarrow u, (n \rightarrow \infty)$となる極限$u \in X$が$M$にも属することを示せれば，定義\ref{def:App:OpenSet}より$M$が閉集合であることが示せる．
点列$(u_n)$は$X$のCauchy列でもあり，$X$がBanach空間であることから完備性より$X$内に極限$v \in X$を持つ．
また，同様に，点列$(u_n)$は$M$のCauchy列でもあること，および定理の仮定から$M$がBanach空間であることから完備性より$M$内に極限$u \in M$を持つ．
そのうえで，$M \subset X$であることと，極限の一意性(定理\ref{thm:limit_unique})より$u = v \in M \subset X$となり，$M$の任意の点列を$(u_n)$の極限$u$は$M$に属する．


\qed
\end{Proof}


%----------------------------------------------------------------------------------------------------------------------------------------------
%----------------------------------------------------------------------------------------------------------------------------------------------
%----------------------------------------------------------------------------------------------------------------------------------------------
\newpage 
\section{作用素の基礎}
本章では作用素というものを導入します．
作用素とは写像の一般化にあたる概念です．
そのため，行列も作用素に含まれますので，ピンと来ない場合は行列として考えてみるといいかもしれません．
ただし，写像の一般化であるため，写像と作用素の違いに注意しながら読むことをお勧めします．
本章の目標の一つは精度保証付き数値計算の証明で良く利用されるNeumann級数に関する定理を紹介することです．


%----------------------------------------------------------------------------------------------------------------------------------------------
\subsection{作用素とは?写像との違いに注意しながら...}
まず，作用素を定義しましょう:

\begin{defi}[作用素]\label{def:operator}
ある線形空間$X$から別の線形空間$Y$への作用素$A$とは，
\begin{align*}
{\mathcal D}(A):=\{ u \in X ~ | ~ Au \in Y \}
\end{align*}
としたとき，${\mathcal D}(A)$のどんな元に対しても，それぞれ集合$Y$のただ一つの元を指定する規則のことである．
また，${\mathcal D}(A)$は$A$の定義域と呼ばれ
\begin{align*}
{\mathcal R}(A):=\{ Au \in Y ~ | ~ u \in {\mathcal D}(A) \}
\end{align*}
を値域と呼ぶ．
\end{defi}


作用素は一見すると写像と同じものに感じますが，作用素は写像の一般化であることに注意が必要です．
実際に，${\mathcal D}(A) \subsetneq X$の場合，上記の作用素$A$は$X$から$Y$への写像ではありません．
また，${\mathcal D}(A) = X$の場合にのみ，作用素$A$は$X$から$Y$への写像となります．
そのため，写像と作用素を同じものと思って取り扱うと引っかかってしまうことが多々あるので注意が必要です．
写像と作用素の違う例の一つが逆作用素です．
写像の場合は単射かつ全射ならば逆写像を持ちますが，作用素の場合は定義が変わってしまいます．

まず，作用素における単射と全射の定義を見てみましょう:

\begin{defi}[単射]\label{def:inj_operator}
線形空間$X$から線形空間$Y$への作用素$A$が
\begin{align*}
u_1 \neq u_2, ~ \forall u_1, u_2 \in {\mathcal D}(A) \Rightarrow A( u_1 ) \neq A( u_2 )
\end{align*}
を満たすときに作用素$A$は単射であるという．
\end{defi}

\begin{defi}[全射]\label{def:sur_operator}
線形空間$X$から線形空間$Y$への作用素$A$が
\begin{align*}
Y = {\mathcal R}(A)
\end{align*}
を満たすときに作用素$A$は全射であるという．
\end{defi}

単射と全射の定義は作用素の特性である$X$と定義域${\mathcal D}(A)$が違うこと以外は写像と同じです．
それに対し，逆作用素の定義は次のようになります:

\begin{defi}[逆作用素]\label{def:inv_operator}
線形空間$X$から線形空間$Y$への作用素$A$とし，その定義域を${\mathcal D}(A) \subset X$，値域を${\mathcal R}(A) \subset Y$とする．
そのとき，
\begin{align*}
A^{-1} (A (u)) = u , ~ u \in {\mathcal D}( A ) \\
A (A^{-1}  (v)) = v , ~ v \in {\mathcal R}( A )
\end{align*}
かつ
\begin{align*}
  {\mathcal D}(A^{-1}) = {\mathcal R}(A) \\
  {\mathcal R}(A^{-1}) = {\mathcal D}(A)
\end{align*}
となる$Y$から$X$への作用素$A^{-1}$を$A$の逆作用素と呼ぶ．
\end{defi}

この逆作用素の定義のもとで，次の定理をみると逆作用素と逆写像の違いが明確になると思います:

\begin{Thm}[単射と逆作用素の関係]\label{thm:inj_inve}
  線形空間$X$から線形空間$Y$への作用素$A$とすると，
\begin{align*}
  A\mbox{が逆作用素を持つ} ~ \Leftrightarrow ~ A\mbox{が単射である}
\end{align*}
\end{Thm}

\begin{Proof}
「$A$が逆作用素を持つ $\Rightarrow$ $A$が単射である」の証明\\
単射の定義\ref{def:inj_operator}の対偶「任意の$u_1, u_2 \in {\mathcal D}(A)$に対し$A (u_1) = A (u_2) \Rightarrow u_1 = u_2$」を満たすことを確かめる．
$A$の逆作用素を$A^{-1}$とすると，任意の$u_1, u_2 \in {\mathcal D}(A)$に対し
\begin{align*}
 & A (u_1) = A (u_2) \\
 \Rightarrow& A^{-1} (A ( u_1 ) ) = A^{-1} (A (u_2)) \\
 \Rightarrow& u_1 = u_2
\end{align*}
となる．
最後の変形は逆作用素$A^{-1}$の定義に由来する．\\


\noindent 「$A$が単射である $\Rightarrow$ $A$が逆作用素$A^{-1}$を持つ」の証明\\
$A$の値域の定義${\mathcal R}(A)=\{ A(u) \in Y ~ | ~ u \in {\mathcal D}(A) \}$より，任意の$v \in {\mathcal R}(A)$に対し
\begin{align*}
  A(u) = v
\end{align*}
となる$u \in {\mathcal D}(A)$が存在する．
その上，$A$が単射であるため，単射の定義\ref{def:inj_operator}の対偶「任意の$u_1, u_2 \in {\mathcal D}(A)$に対し$A (u_1) = A (u_2) \Rightarrow u_1 = u_2$」より，$u \in {\mathcal D}(A)$はどんな$v \in {\mathcal R}(A)$に対してもただ一つの元である．
そのため，作用素の定義\ref{def:operator}より，上記の$v \in {\mathcal R}(A)$に対してただ一つの元$u \in {\mathcal D}(A)$を指定する規則として
\begin{align*}
  B(v) = u
\end{align*}
となる定義域${\mathcal D}(B) = {\mathcal R}(A)$と値域${\mathcal R}(B)={\mathcal D}(A)$となる$Y$から$X$への作用素$B$が定義できる．
その上，$B(v) = u$の$v$に$v = A(u)$を代入すると
\begin{align*}
  B(A(u)) = u
\end{align*}
となる．
同様に，$A(u) = v$の$u$に$u = B(v)$を代入すると
\begin{align*}
  A( B(v) ) = v
\end{align*}
となる．
よって，定義域${\mathcal D}(B) = {\mathcal R}(A)$と値域${\mathcal R}(B)={\mathcal D}(A)$となる$Y$から$X$への作用素$B$は$A$の逆作用素であるため，$A$は逆作用素を持つ．

\qed
\end{Proof}

写像の場合，逆写像を持つ必要十分条件は単射かつ全射でした．
それに対し，逆作用素を持つ必要十分条件には作用素$A$が全射であることを必要としません．
なぜなら，逆作用素$A^{-1}$は作用素であるため定義域${\mathcal D}(A^{-1})$と$Y$が一致しなくてもよいからです．
これは，写像と作用素が明確に違うことを表しています．
そのために，1から作用素を扱うには，基礎から丁寧に定義しなければなりません．

では，なぜ，苦労してまで写像を捨てて，$X$と${\mathcal D}(A)$が別になる作用素を利用するのか疑問に思うと思います．
もちろん，苦労する分，うまみが出てくるから作用素を導入します．
$X$と${\mathcal D}(A)$を別にする理由は固有値やその一般化であるスペクトルにあります．
実際に，線形代数で習った行列の固有値問題は行列$A$は正方行列であったと思います．
正方行列ということは行列$A$は$X = {\mathcal D}(A) = {\mathbb R}^N$から$Y = {\mathbb R}^N$への作用素になります．
そのため，線形代数で習った固有値問題は$X = {\mathcal D}(A) = Y$となる状況です．
しかし，単純な2階微分作用素$d^2/dx^2$を考えてみると，${\mathcal D}(A)$には2階以上微分できる関数である必要があるのに対し，$Y$の元は微分ができる必要はありません．
では，このような2階微分作用素$d^2/dx^2$に対して固有値問題を考えたい場合，どうすればよいでしょうか?
写像のままで進めようとすると$X = {\mathcal D}(A) \neq Y$で写像$A$の固有値を定義する必要があります．
これがまかり通ると，正方行列以外でも固有値が定義できていることになってしまい，線形代数の固有値問題とは別ものになってしまいます．
このような状況で作用素のうまみがでてきます．
$X=Y$を連続関数全体にし，${\mathcal D}(A)$のみ2階微分可能な関数空間に設定すれば，今までの線形代数の固有値問題を壊さずに定義していくことができます．
行列の場合は，たまたま$X = {\mathcal D}(A)$になったという解釈です．

では，次に作用素$A$と$B$が等しい(すなわち$A=B$)とはどういうことか，ちゃんと定義をしましょう．

\begin{defi}[作用素の等号]\label{def:operator_equal}
線形空間$X$から線形空間$Y$への作用素$A$と$B$が等しいとは
\begin{align*}
{\mathcal D}(A) = {\mathcal D}(B)
\end{align*}
かつ
\begin{align*}
A u = B u, ~ \forall u \in {\mathcal D}(A) = {\mathcal D}(B)
\end{align*}
が成立することであり，
\begin{align*}
A = B
\end{align*}
と表記する．
\end{defi}

定義をみてもわかるとおり，定義域も一致していることが重要です．

この節の最後に作用素の連続性について定義します．
作用素の連続性については，本書籍ではあまり利用しないので簡単に留めておく程度で十分です．

\begin{defi}[作用素の連続]\label{def:operator_continue}
ノルム空間$X$からノルム空間$Y$への作用素$A$が$u \in {\mathcal D}(A)$で連続であるとは
\begin{align*}
	u_n \rightarrow u, ~ (n \rightarrow \infty)
\end{align*}
となる任意の$u_n \in {\mathcal D}(A) \subset X$に対して
\begin{align*}
	A u_n \rightarrow A u, ~ (n \rightarrow \infty)
\end{align*}
を満たすときである．
さらに，$A$が任意の$u \in {\mathcal D}(A)$において連続であるとき，$A$は連続であるという．
\end{defi}

上記の作用素$A$の$u \in {\mathcal D}(A)$における連続と同値になる$\varepsilon - \delta$論法を用いた定義は次のようになります:
\begin{align*}
\forall \varepsilon > 0, ~ \exists \delta > 0, ~ \| u_n - u \|_{X} < \delta ~ \mbox{となる} ~ \forall u_n \in {\mathcal D}(A) ~ \mbox{に対して} ~ \| A u_n - A u \|_{Y} < \varepsilon
\end{align*}

「作用素$A$は$u \in {\mathcal D}(A)$で連続」と「作用素$A$は連続」は連続である点が，$u \in {\mathcal D}(A)$のただ1点における連続と，定義域${\mathcal D}(A)$の任意の元で連続という大きな違いがあるので注意して下さい．

連続な作用素も作用素の一つのクラスになります．
作用素はあまりにも広すぎる定義であるため，クラス分けして，「このクラスの作用素ならば，こんな性質を持ちます」といった感じで解析していきます．


%----------------------------------------------------------------------------------------------------------------------------------------------
\subsection{線形作用素}
本節では最も標準的な作用素のクラスである線形作用素を紹介します．
線形作用素では作用素同士の加法とスカラー乗法なども定義できます．
では，線形作用素の定義を見てみましょう:

\begin{defi}[線形作用素]\label{def:linearoperator}
線形空間$X$から線形空間$Y$への作用素$A$が，任意の$u, v \in {\mathcal D}(A) \subset X$と$\alpha \in \mathbb{K}$に対し，
\begin{align*}
& {\mathcal D}(A) ~ \mbox{が} X \mbox{の線形部分空間}\\
& A(u + v) = A u + A v \\
& A(\alpha u) = \alpha A u
\end{align*}
を満たすとき，$A$を線形作用素と呼ぶ．
\end{defi}

線形作用素の定義は線形写像の定義を作用素に置き換えたものです．
もちろん${\mathcal D}(A) \subset X$であることに注意してください．
同じように線形作用素に対する加法，スカラー乗法，合成作用素を定義していきましょう．
これらを定義する際は，特に定義域に注意しましょう．

\begin{defi}[線形作用素の加法]\label{def:add_operator}
	線形空間$X$から線形空間$Y$への線形作用素$A$と$B$の和を
	\begin{align*}
	(A + B) u := A u + B u, ~ u \in {\mathcal D}(A) \cap {\mathcal D}(B)
	\end{align*}
	と定義する．
	このとき，$X$から$Y$への作用素$A+B$の定義域は
	\begin{align*}
	{\mathcal D}(A+B) =  {\mathcal D}(A) \cap {\mathcal D}(B)
	\end{align*}
	とする．
\end{defi}


\begin{defi}[線形作用素のスカラー乗法]\label{def:scalar_operator}
	線形空間$X$から線形空間$Y$への線形作用素$A$の$\alpha \in {\mathbb K}$によるスカラー倍を
	\begin{align*}
	(\alpha A) u := \alpha (Au), ~ u \in {\mathcal D}(A)
	\end{align*}
	と定義する．
	このとき，$X$から$Y$への作用素$\alpha A$の定義域は
	\begin{align*}
	{\mathcal D}(\alpha A) :=  {\mathcal D}(A)
	\end{align*}
	とする．
\end{defi}


\begin{defi}[合成作用素]\label{def:mul_operator}
	$X,Y,Z$を線形空間とする．
	$A$を$Y$から$Z$への線形作用素とし，$B$を$X$から$Y$への線形作用素とする．
	そのとき，$A$と$B$の合成作用素$AB$は
	\begin{align*}
	(A B) u := A( B u), ~ u \in \{ v \in {\mathcal D}(B) ~ | ~ Bv \in {\mathcal D}(A)  \}
	\end{align*}
	と定義する．
	このとき，$X$から$Z$への合成作用素$AB$の定義域は
	\begin{align*}
	{\mathcal D}(AB) := \{ v \in  {\mathcal D}(B) ~ | ~ Bv \in {\mathcal D}(A)  \}
	\end{align*}
	とする．
\end{defi}

線形作用素同士の合成作用素は線形作用素になるとは限らないので注意してください．
なぜなら，${\mathcal D}(B)$が$X$の線形部分空間でも，${\mathcal D}(AB)$が$X$の線形部分空間になるとは限らないからです．
${\mathcal D}(AB)$が$X$の線形部分空間であれば，合成作用素$AB$は線形作用素になります．


線形作用素の場合，単射性をチェックするための定理がいくつかあります．
線形作用素が単射であることを確かめる際には，単射の定義をチェックするよりも楽な場合が多いので紹介します:

\begin{Thm}[線形作用素に対する単射性(1)]\label{thm:linear_operator_injective11}
	線形空間$X$から線形空間$Y$への線形作用素$A$において以下は同値である:\\
	i) 線形作用素$A$が単射である \\
	ii) $A u = 0, ~ u \in {\mathcal D}(A)  ~ \Rightarrow ~ u = 0$
\end{Thm}

\begin{Proof}
	単射の定義の対偶は
	\begin{align*}
	A u_1  = A u_2 , ~ \forall u_1, u_2 \in {\mathcal D}(A) \Rightarrow  u_1 = u_2
	\end{align*}
	となる．
	その上，$A$は線形作用素であるため
	\begin{align*}
	A u_1  = A u_2 ~ \Leftrightarrow ~ A (u_1 - u_2) = 0
	\end{align*}
	となる．
	$u_1 - u_2 \in {\mathcal D}(A)$を$u \in {\mathcal D}(A)$とおきなおせば，(i) $\Rightarrow$ (ii)は証明された．
	また，証明を逆に追うことで，(ii) $\Rightarrow$ (i) も示せる．

\qed
\end{Proof}

\begin{Thm}[線形作用素に対する単射性(2)]\label{thm:linear_operator_injective2}
	ノルム空間$X$からノルム空間$Y$への線形作用素$A$とする．
	不等式
	\begin{align*}
	 \| u \|_{X} \le K \| A u \|_{Y}, ~ u \in {\mathcal D}(A)
	\end{align*}
	を満たす定数$K > 0$が存在するならば，線形作用素$A$は単射である．
\end{Thm}

\begin{Proof}
	$A$が線形作用素であるため，定理\ref{thm:linear_operator_injective11} ii)を使って証明する．
	ノルムの定義より
	\begin{align*}
	 A u = 0, ~ \forall u \in {\mathcal D}(A) \Leftrightarrow \| A u \|_{Y} = 0
	\end{align*}
	となる．
	さらに，$A u = 0$ならば，
	\begin{align*}
	 \| u \|_{X} \le K \| A u \|_{Y} = 0, ~ u \in {\mathcal D}(A)
	\end{align*}
	より$\| u \|_{X} = 0$となる．
	よって，再びノルムの定義より
	\begin{align*}
	 \| u \|_{X} = 0, ~ \forall u \in {\mathcal D}(A) \Leftrightarrow u  = 0
	\end{align*}
	より，$A u = 0$ならば$u = 0$となる．

\qed
\end{Proof}

定理\ref{thm:linear_operator_injective2}と逆作用素の定義より，線形作用素$A$に対して不等式
\begin{align*}
 \| u \|_{X} \le K \| A u \|_{Y}, ~ u \in {\mathcal D}(A)
\end{align*}
を満たす定数$K > 0$が存在するならば，線形作用素$A$は$Y$から$X$への逆作用素$A^{-1}$を持ちます．
但し，$A^{-1}$の定義域${\mathcal D}(A^{-1}) = {\mathcal R}(A)$は$Y$と一致するかどうかはわかりません．
もし，$A$が全射でもあるならば，もちろん$A^{-1}$の定義域${\mathcal D}(A^{-1}) = {\mathcal R}(A)$と$Y$は一致します．




%----------------------------------------------------------------------------------------------------------------------------------------------
\subsection{有界な線形作用素}
前節では作用素に「線形」という条件を追加しました．
本節では，作用素に「線形」かつ「有界」という条件を追加した場合の作用素のクラスの性質を紹介します．
では，まず，有界な線形作用素を定義しましょう:

\begin{defi}[有界な線形作用素]\label{def:bounded_operator}
	ノルム空間$X$から$Y$への線形作用素$A$に対し，
	\begin{align*}
	\| A u \|_{Y} \le K \| u \|_{X}, ~ u \in {\mathcal D}(A)
	\end{align*}
	を満たす正の定数$K$が存在するとき，線形作用素$A$を有界な作用素と呼ぶ．
\end{defi}

有界な線形作用素の定義は非常に悩ましく，$X = {\mathcal D}(A)$の場合にのみ限り，有界な線形作用素と呼ぶ文献も数多くあります．
もちろん$X = {\mathcal D}(A)$に限らない定義の文献もあります．
しかし，本書籍では$X \neq {\mathcal D}(A)$でも次の定理が成り立つことから，有界な線形作用素の定義には$X = {\mathcal D}(A)$は入れないので注意しでください．

\begin{Thm}[有界な線形作用素と連続な線形作用素]\label{thm:bounded_continus_operator}
	ノルム空間$X$からノルム空間$Y$への線形作用素$A$に対し，
	\begin{align*}
	A \mbox{が有界} ~ \Leftrightarrow ~ A \mbox{が連続}
	\end{align*}
\end{Thm}

\begin{Proof}
	「$A$が有界$\Rightarrow A$が連続」の証明\\
	連続性の定義より，$u_n \rightarrow u$となる任意の$u_n \in {\mathcal D}(A)$に対して$A u_n \rightarrow A u$となることを確かめる．
	$u_n \rightarrow u$となる任意の$u_n \in {\mathcal D}(A)$から$\| u_n - u \|_{X} \rightarrow 0, ~ (n \rightarrow \infty)$を持つ．
	その上，$A$は有界であることから
	\begin{align*}
	 \| A u_n - A u \|_{Y} \le M \| u_n - u \|_{X} \rightarrow 0, ~ (n \rightarrow \infty)
	\end{align*}
	となる．
	よって，$u_n \rightarrow u, ~ (n \rightarrow \infty)$ならば，$A u_n \rightarrow A u$であるため，$A$は連続である．\\
	
	\noindent 「$A$が連続$\Rightarrow A$が有界」の証明\\
	背理法によって証明する．
	すなわち，$A$が連続ならば，任意の$M_2 > 0$に対して
	\begin{align*}
	 \| A u \|_{Y} > M_2 \| u \|_{X} 
	\end{align*}
	を満たす$u \in {\mathcal D}(A)$が存在すると仮定して矛盾をみつける．
	この仮定より自然数$n$に対して，
	\begin{align*}
	 \| A u_n \|_{Y} > n \| u_n \|_{X} 
	\end{align*}
	を満たす$u_n \in {\mathcal D}(A)$が存在する．
	このとき，$\| u_n \|_X \neq 0$であることに注意する．
	ノルム空間$X$はノルム空間の定義より線形空間であるため，ゼロ元$0 \in X$を持つ．
	その上，線形作用素の定義より${\mathcal D}(A)$は$X$の部分空間であるため，ゼロ元$0 \in {\mathcal D}(A) \subset X$を持つ．
	その上，$A$が連続であるため，$A$は$0 \in {\mathcal D}(A)$でも連続である．
	$\varepsilon - \delta$論法による$A$の$0 \in {\mathcal D}(A) \subset X$における連続の定義を記述すると
	\begin{align*}
	\forall \varepsilon > 0, ~ \exists \delta > 0, ~ \| u_n \|_{X} < \delta ~ \mbox{となる} ~ \forall u_n \in X ~ \mbox{に対して} ~ \| A u_n \|_{Y} < \varepsilon
	\end{align*}
	となる．
	その上，$\varepsilon$を$n \| u_n \|_{X} $とすると，$\delta_n > 0$が存在し，$\| u_n \|_{X} < \delta$となる任意の$u_n \in {\mathcal D}(A)$に対して，
	\begin{align*}
	 \| A u_n \|_{Y} < n \| u_n \|_{X}
	\end{align*}
	となる．
	有界ではないという仮定と組み合わせると
	\begin{align*}
	n \| u_n \|_{X} < \| A u_n \|_{Y} < n \| u_n \|_{X}
	\end{align*}
	となるため矛盾する．

\qed
\end{Proof}



%----------------------------------------------------------------------------------------------------------------------------------------------
\subsection{定義域が\texorpdfstring{$X$}{X}の全体となる有界な線形作用素全体の集合\texorpdfstring{${\mathcal B}(X, Y)$}{B(X,Y)}}\label{sec:DX_bounded_linear_op}
今まで作用素に「連続」，「線形」，「有界かつ線形」などのクラス分けをしました．
この節では，さらに「有界かつ線形かつ$X = {\mathcal D}(A)$」となる作用素を考えます．
しかし，今までと大きく違う点は，「有界かつ線形かつ$X = {\mathcal D}(A)$」となる作用素となるものすべて集めた集合を考えるので注意してください．
では，まず，本節で主役となる集合の定義をしましょう．

\begin{defi}[定義域が$X$の全体となる有界な線形作用素全体の集合${\mathcal B}(X, Y)$]\label{def:set_bound_linear_operator}
	Banach空間$X$からBanach空間$Y$の有界な線形作用素$A$の定義域が$X$，すなわち
	\begin{align*}
	 X = {\mathcal D}(A)
	\end{align*}
	となるもの全体の集合を
	\begin{align*}
	 {\mathcal B}(X, Y)
	\end{align*}
	と書く．
	また，$X = Y$の場合は省略して${\mathcal B}(X)$と記述する．
\end{defi}

集合${\mathcal B}(X)$の元$A$は，$X = {\mathcal D}(A)$となる有界な線形作用素であるため写像でもあります．
しかし，写像と呼ぶことはなく，あくまで作用素と呼びます．
その理由は，$X = {\mathcal D}(A)$となる有界な線形作用素$A$が全射ではないが単射の場合は，逆作用素$A^{-1}$を持ちますが，写像$A$と呼んでしまうと逆写像$A^{-1}$は持ちません．
このように，$X = {\mathcal D}(A)$となる有界な線形作用素と写像は一見すると同じに見えますが，やはり$X = {\mathcal D}(A)$となる有界な線形作用素も写像のより一般的な定義になっていることに注意が必要です．


本節の主役である集合${\mathcal B}(X, Y)$に着目する理由は，${\mathcal B}(X, Y)$に適切なノルムを入れるとBanach空間になるからです．
ここでは，まず，${\mathcal B}(X, Y)$に適切なノルムを入れ，ノルム空間になることを示した上で，Banach空間になることを証明します．
\ref{sec:Banach_space}節の山場であるため，最初は証明を追うことができないかもしれません．
証明を飛ばして結果だけ見る場合でも，証明後に掲載しているNeumann級数に関する定理は数値計算の品質保証で度々でてくるため，必ず把握しておいて下さい．

\begin{Thm}[${\mathcal B}(X, Y)$はBanach空間]\label{thm:set_bound_linear_operator}
	$X$をノルム空間とし，$Y$をBanach空間とする．
	定義域が$X$全体となる$X$から$Y$への有界な線形作用素全体の集合${\mathcal B}(X, Y)$のノルムを
	\begin{equation}
		\| A \|_{{\mathcal B}(X,Y)} := \sup_{ u \in X \setminus \{ 0 \}} \frac{ \| Au \|_{Y} }{ \| u \|_{X} }, ~ A \in {\mathcal B}(X, Y) \label{def:operator_norm}
	\end{equation}
	とすると${\mathcal B}(X, Y)$はBanach空間となる．
\end{Thm}

\begin{Proof}

定義\ref{def:add_operator}の作用素の加法と定義\ref{def:scalar_operator}の作用素のスカラー乗法をもとに線形空間の公理\ref{axiom:linear_space}が満たされていることは簡単に導かれる．
但し，${\mathcal B}(X, Y)$のゼロ元は任意の$u \in X$を$0 \in Y$へ写す作用素であることに注意が必要である．

\noindent 「ノルム空間」\\
$\| A \|_{{\mathcal B}(X,Y)}$がノルムの定義\ref{def:App:norm}を満たすことを示せばよい．
ノルム空間$X$とBanach空間$Y$であるため$\| \cdot \|_{X} \ge 0$と$\| \cdot \|_{Y} \ge 0$であることから
\begin{align*}
 \frac{ \| Au \|_{Y} }{ \| u \|_{X} } \ge 0
\end{align*}
となるため，$\| A \|_{{\mathcal B}(X,Y)} \ge 0$となり，ノルムの定義\ref{def:App:norm} (i)はいえる．

次に，$A = 0$ならば$\| A u \|_{Y} = 0$であるため，
\begin{align*}
\| A \|_{{\mathcal B}(X,Y)} = \sup_{ u \in X \setminus \{ 0 \}} \frac{ \| Au \|_{Y} }{ \| u \|_{X} } = \sup_{ u \in X \setminus \{ 0 \}} \frac{ 0 }{ \| u \|_{X} } =  0
\end{align*}
である．さらに，
任意の$u \in X \setminus \{ 0 \}$について
\begin{align*}
\frac{ \| Au \|_{Y} }{ \| u \|_{X} } = 0 \Leftrightarrow \| Au \|_{Y} = 0 \Leftrightarrow Au = 0
\end{align*}
任意の$u \in X \setminus \{ 0 \}$を$0 \in Y$へ写す作用素は${\mathcal B}(X, Y)$が線形空間より一意に存在し，$A = 0$である．
よって，ノルムの定義\ref{def:App:norm} (ii)も示された．

続いて，$\alpha \in {\mathbb K}$としたとき，$Y$はBanach空間であるため$\| \cdot \|_Y$はノルムの定義を満たすため，
\begin{align*}
\|  \alpha A \|_{{\mathcal B}(X,Y)} = \sup_{ u \in X \setminus \{ 0 \}} \frac{ \| \alpha Au \|_{Y} }{ \| u \|_{X} } = | \alpha | \sup_{ u \in X \setminus \{ 0 \}} \frac{ \| Au \|_{Y} }{ \| u \|_{X} } = | \alpha | \|   A \|_{{\mathcal B}(X,Y)}
\end{align*}
となるため，ノルムの定義\ref{def:App:norm} (iii)も示された．

最後に任意の$A, B \in {\mathcal B}(X, Y)$について
\begin{align*}
\|  A + B \|_{{\mathcal B}(X,Y)} &= \sup_{ u \in X \setminus \{ 0 \}} \frac{ \| (A + B) u \|_{Y} }{ \| u \|_{X} } = \sup_{ u \in X \setminus \{ 0 \}} \frac{ \| Au + Bu \|_{Y} }{ \| u \|_{X} } \\
&\le \sup_{ u \in X \setminus \{ 0 \}} \frac{ \| Au \|_{Y} + \| Bu \|_{Y} }{ \| u \|_{X} } \\
&\le \sup_{ u \in X \setminus \{ 0 \}} \frac{ \| Au \|_{Y} }{ \| u \|_{X} } + \sup_{ u \in X \setminus \{ 0 \}} \frac{ \| Bu \|_{Y} }{ \| u \|_{X} } \\
&= \|  A \|_{{\mathcal B}(X,Y)} + \|  B \|_{{\mathcal B}(X,Y)}
\end{align*}
となり，ノルムの定義\ref{def:App:norm} (iv)も示されたため，${\mathcal B}(X, Y)$はノルム空間である．\\

\noindent 「Banach空間」\\
Banach空間であることを証明するには${\mathcal B}(X, Y)$の任意のCauchy列$( A_n ) \subset {\mathcal B}(X, Y)$が極限$T$を${\mathcal B}(X, Y)$内に持つことを示せばよい．
証明の流れは，
まず，極限の候補$\tilde{A}$の定義できるか確認する．
続いて極限の候補$\tilde{A}$が${\mathcal B}(X, Y)$に属していることを示す．
最後に，極限の候補$\tilde{A}$がCauchy列$( A_n )$の極限であることを示す．

まず，極限の候補$\tilde{A}$が定義に取り掛かる．
任意のCauchy列$( A_n ) \subset {\mathcal B}(X, Y)$はCauchy列の定義\ref{def:App:Cau}より
\begin{align*}
\|  A_n  - A_m  \|_{{\mathcal B}(X,Y)} \rightarrow 0, ~ (n, m \rightarrow \infty)
\end{align*}
となる．
任意の$u \in X \setminus\{ 0 \}$に対して，点列$(A_n u) \subset Y$は
\begin{align*}
\|  A_n u  - A_m u \|_{Y} &= \frac{ \|  (A_n  - A_m) u \|_{Y} }{ \| u \|_{X} } \| u \|_{X} \\
&\le \sup_{\phi \in X \setminus \{ 0 \}} \frac{ \|  (A_n  - A_m) \phi \|_{Y} }{ \| \phi \|_{X} } \| u \|_X  \\
&= \| A_n - A_m \|_{{\mathcal B}(X,Y)} \| u \|_{X} \rightarrow 0, ~ (n, m \rightarrow \infty)
\end{align*}
を持つため，点列$(A_n u) \subset Y$はCauchy列になる．
その上，$Y$はBanach空間であるため，$Y$の任意のCauchy列は収束し，$Y$内に極限$\tilde{A} u$となるような$X$から$Y$への作用素$\tilde{A}$が存在する．
ここで，任意の$u \in X$に対して極限$\tilde{A} u$が定義されることから，$\tilde{A}$の定義域は${\mathcal D}(\tilde{A}) = X$である．
これで，${\mathcal B}(X, Y)$の任意のCauchy列$( A_n )$の極限の候補$\tilde{A}$が定義できた．

続いて，定義した極限の候補$\tilde{A}$が${\mathcal B}(X, Y)$に属しているか確認する．
$\tilde{A}$が有界な線形作用素であり，かつ${\mathcal D}(\tilde{A}) = X$であることを示せばよい．
${\mathcal B}(X, Y)$の任意のCauchy列$( A_n )$の元$A_n$は線形作用素であるため，線形作用素の定義\ref{def:linearoperator}より任意の$\alpha, \beta \in {\mathbb K}$と$u, v \in X$について
\begin{align*}
A_n (\alpha u + \beta v) = \alpha A_n u + \beta A_n v
\end{align*}
を持つ．よって$n \rightarrow \infty$とすると
\begin{align*}
\tilde{A} (\alpha u + \beta v) = \alpha \tilde{A} u + \beta \tilde{A} v
\end{align*}
となり，極限の候補$\tilde{A}$は線形作用素である．
次に極限の候補$\tilde{A}$が有界作用素であることを示す．
点列$( A_n )$はCauchy列であるため定理\ref{def:if_cauchy_then_bounded}より有界列でもある．
すなわち，どんな$n \in {\mathbb N}$に対しても
\begin{align*}
\| A_n \|_{{\mathcal B}(X,Y)} \le M
\end{align*}
となる$n \in {\mathbb N}$に依存しない定数$M$が存在する．
この$n \in {\mathbb N}$に依存しない定数$M$は，任意の$u \in X$について
\begin{align*}
\| A_n u \|_{Y} \le M \| u \|_{X}
\end{align*}
も満たす．
$A_n u \rightarrow \tilde{A} u, ~ (n \rightarrow \infty)$であるため，上の不等式に対して$n \rightarrow \infty$とすると$M$が$n$に依存しないため
\begin{align*}
\| \tilde{A} u \|_{Y} \le M \| u \|_{X}
\end{align*}
を得る．
よって，点列$(A_n)$の極限の候補$\tilde{A}$は${\mathcal B}(X, Y)$に属する．

最後に，点列$(A_n)$の極限が$\tilde{A}$であることを示す．
任意の$u \in X$に対して，点列$(A_n u) \subset Y$は$Y$内に極限$\tilde{A} u$を持つこと，すなわち
\begin{align*}
A_n u \rightarrow \tilde{A} u, ~ (n \rightarrow \infty)
\end{align*}
を持つことから
\begin{align*}
 \| A_n u - A_m u \|_{Y} \rightarrow \| A_n u - \tilde{A} u \|_{Y}, ~ (m \rightarrow \infty)
\end{align*}
となる．
その上，${\mathcal B}(X,Y)$のノルムの定義と$\tilde{A} \in {\mathcal B}(X, Y)$から
\begin{align*}
\| A_n - A_m \|_{{\mathcal B}(X,Y)}  \rightarrow \| A_n - \tilde{A} \|_{{\mathcal B}(X,Y)}, ~ (m \rightarrow \infty)
\end{align*}
を得る．
点列$(A_n)$がCauchy列であるため
\begin{align*}
  \forall \varepsilon > 0, ~ \exists N \in {\mathbb N}, ~ \forall n, m \ge N ~ \mbox{に対して} ~ \| A_n - A_m \|_{{\mathcal B}(X,Y)} < \varepsilon
\end{align*}
を満たす．
その上，$m \rightarrow \infty$とすると
\begin{align*}
  \forall \varepsilon > 0, ~ \exists N \in {\mathbb N}, ~ \forall n \ge N ~ \mbox{に対して} ~ \| A_n - \tilde{A} \|_{{\mathcal B}(X,Y)} < \varepsilon
\end{align*}
となり，$\tilde{A} \in {\mathcal B}(X,Y)$はCauchy列$(A_n)$の極限である．
よって，任意のCauchy列は${\mathcal B}(X,Y)$内に極限を持つため，ノルム空間${\mathcal B}(X,Y)$はBanach空間である．

\qed
\end{Proof}

これで，\ref{sec:Banach_space}節の山場は越えました．
続いて，Banach空間${\mathcal B}(X,Y)$のノルムに関するよく使われる性質を2つほど紹介しておきます．
紹介する2つの性質はほぼ当たり前なので，証明は冗長に感じるかもしれません．

\begin{Thm}[${\mathcal B}(X, Y)$のノルムの性質(1)]\label{thm:linear_operator_norm_prop1}
$X$をノルム空間とし，$Y$をBanach空間とする．
そのとき，任意の$u \in X$と任意の$A \in {\mathcal B}(X, Y)$について以下の不等式が成り立つ:
\begin{align*}
\| A u \|_{Y} \le \| A \|_{{\mathcal B}(X,Y)} \| u \|_{X}
\end{align*}
\end{Thm}

\begin{Proof}
$u = 0$の場合は明らかに成り立つため，$u \in X \setminus \{ 0 \}$について考える．
$u \in X \setminus \{ 0 \}$について
\begin{align*}
\| A u \|_{Y} = \frac{ \| A u \|_{Y} }{ \| u \|_{X} } \| u \|_{X} \le \sup_{\phi \in X \setminus \{ 0 \}} \frac{ \| A \phi \|_{Y} }{ \| \phi \|_{X} } \| u \|_{X} = \| A \|_{{\mathcal B}(X,Y)} \| u \|_{X}
\end{align*}
となるため，題意は示された．

\qed
\end{Proof}

\begin{Thm}[${\mathcal B}(X, Y)$のノルムの性質(2)]\label{thm:linear_operator_norm_prop2}
$X$をノルム空間とし，$Y$と$Z$をBanach空間とする．
そのとき，任意の$B \in {\mathcal B}(X, Y)$と$A \in {\mathcal B}(Y, Z)$の合成作用素$AB$は${\mathcal B}(X, Z)$に属する．
その上，
\begin{align*}
\| A B \|_{{\mathcal B}(X, Z)} \le \| A \|_{{\mathcal B}(Y, Z)} \| B \|_{{\mathcal B}(X,Y)}
\end{align*}
\end{Thm}

\begin{Proof}
合成作用素の定義\ref{def:mul_operator}から
\begin{align*}
{\mathcal D}(AB) = \{ v \in {\mathcal D}(B) = X ~ | ~ Bv \in {\mathcal D}(A) = Y \}
\end{align*}
となるが，$B \in {\mathcal B}(X, Y)$であるため，任意の$v \in X$に対して$B v$は$Y$に属する．
よって，
\begin{align*}
{\mathcal D}(AB) = {\mathcal D}(B) = X
\end{align*}
となる．
その上，$A$も$B$も線形作用素であることから，任意の$u, v \in X$と任意の$\alpha, \beta \in {\mathbb K}$に対して
\begin{align*}
AB ( \alpha u + \beta v ) = A ( B \alpha u + B \beta v ) = A ( \alpha B u + \beta B v ) = A \alpha B u + A \beta B v = \alpha A B u + \beta A B v
\end{align*}
となるため，合成作用素$AB$は定義域が$X$全体となる線形作用素である．
また，$A \in {\mathcal B}(Y, Z)$, $B \in {\mathcal B}(X, Y)$であるため，任意の$u \in X$について，定理\ref{thm:linear_operator_norm_prop1}から
\begin{align*}
 \| A B u \|_{Z} \le \| A \|_{{\mathcal B}(Y, Z)} \| B u \|_{Y} \le  \| A \|_{{\mathcal B}(Y, Z)} \| B \|_{{\mathcal B}(X, Y)} \| u \|_{X}
\end{align*}
となり，定義域が$X$全体となる線形作用素$AB$は有界な線形作用素でもある．
よって$AB$は${\mathcal B}(X, Z)$に属する．
その上，
\begin{align*}
\| A B \|_{{\mathcal B}(X, Z)} &= \sup_{ u \in X \setminus \{ 0 \}} \frac{ \| A B u \|_{Z} }{ \| u \|_{X} } \\
&\le \sup_{ u \in X \setminus \{ 0 \}} \frac{ \| A \|_{{\mathcal B}(Y, Z)} \| B \|_{{\mathcal B}(X, Y)} \| u \|_{X} }{ \| u \|_{X} } \\
&= \| A \|_{{\mathcal B}(Y, Z)} \| B \|_{{\mathcal B}(X, Y)}
\end{align*}

\qed
\end{Proof}

それでは，\ref{sec:Banach_space}章の最後にNeumann級数に関する定理を紹介に入ります．
Neumann級数に関する定理を紹介にあたり，まず，恒等作用素を定義します:

\begin{defi}[$X$上の恒等作用素]\label{def:identity_operator_on_X}
$X$をBanach空間とする．
任意の$u \in X$に対して
\begin{align*}
 I u = u
\end{align*}
となる$I \in {\mathcal B}(X)$を$X$上の恒等作用素と呼ぶ．
\end{defi}

もちろん，$X$上の恒等作用素$I \in {\mathcal B}(X)$は，任意の$A \in {\mathcal B}(X)$に対して
\begin{align*}
 A I  = I A = A
\end{align*}
となります．
また，もちろん，
\begin{align*}
 {\mathcal R}(I) = X
\end{align*}
です．

では，Neumann級数に関する定理を紹介します．
Neumann級数に関する定理は数値計算の品質保証で重要になる定理ですので覚えておいて下さい．

\begin{Thm}[Neumann級数]\label{thm:nuemann_ser}
$X$をBanach空間とする．
$B \in {\mathcal B}(X)$とし，$I \in {\mathcal B}(X)$を$X$上の恒等作用素とする．
もし
\begin{align*}
 \| I - B \|_{{\mathcal B}(X)} < 1
\end{align*}
ならば，$B$は逆作用素をもち$B^{-1} \in {\mathcal B}(X)$(すなわち，$B$は全単射)となる．
そのうえ，
\begin{align*}
  B^{-1} = I + (I - B) + (I - B)^2 + \cdots = \sum_{i = 0}^\infty (I - B)^i
\end{align*}
で，かつ
\begin{align*}
  \| B^{-1} \|_{{\mathcal B}(X)} \le \frac{ 1 }{ 1 - \| I - B \|_{{\mathcal B}(X)} }
\end{align*}

\end{Thm}

\begin{Proof}
\begin{align*}
  S_n = I + ( I - B) + ( I - B )^2 + \cdots + ( I - B )^n
\end{align*}
とすると，$B$と$I$はともに${\mathcal B}(X)$に属するため，加法$I - B$や合成作用素$(I - B)(I - B)$なども${\mathcal B}(X)$に属する．
よって$S_n$も${\mathcal B}(X)$に属する．

続いて，点列$(S_n) \subset {\mathcal B}(X)$が極限$S$を${\mathcal B}(X)$内に持つか確認する．
定理\ref{thm:linear_operator_norm_prop2}より
\begin{align*}
  \| (I - B)^i \|_{{\mathcal B}(X)} \le \| I - B \|_{{\mathcal B}(X)}^i, ~ i = 0, 1, \cdots
\end{align*}
となるため，$n > m > 0$となる整数に対して
\begin{align*}
  \| S_n - S_m \|_{{\mathcal B}(X)} = \left\| \sum_{i = m + 1}^n  ( I - B )^i \right\|_{{\mathcal B}(X)} \le \sum_{i = m + 1}^n \| I - B \|_{{\mathcal B}(X)}^i
\end{align*}
となる．
定理の仮定より$ \| I - B \|_{{\mathcal B}(X)} < 1$であるため，
\begin{align*}
  \sum_{i = m + 1}^n \| I - B \|_{{\mathcal B}(X)}^i \rightarrow 0, ~ (n, ~ m \rightarrow \infty )
\end{align*}
となる．
よって，
\begin{align*}
  \| S_n - S_m \|_{{\mathcal B}(X)} \rightarrow 0, ~ (n, ~ m \rightarrow \infty )
\end{align*}
となるため，点列$(S_n)$はCauchy列である．
その上，${\mathcal B}(X)$はBanach空間であるため，任意のCauchy列は極限を${\mathcal B}(X)$に持つため，点列$(S_n)$は
\begin{align*}
  \| S_n - S \|_{{\mathcal B}(X)} \rightarrow 0, ~ (n, ~ m \rightarrow \infty )
\end{align*}
となる極限$S \in {\mathcal B}(X)$を持つ．

次に$S$が$B^{-1}$になることを示す．
合成作用素の定義\ref{def:mul_operator}に従って合成作用素$BS_n$を考える．
$X$はBanach空間であり，$B, S_n \in {\mathcal B}(X)$であるため，定理\ref{thm:linear_operator_norm_prop2}合成作用素$BS_n$は${\mathcal B}(X)$に属する．
その上，点列$(B S_n) \subset {\mathcal B}(X)$は
\begin{align*}
  \| B S_n - B S \|_{{\mathcal B}(X)} \le \| B \|_{{\mathcal B}(X)}  \| S_n - S \|_{{\mathcal B}(X)} \rightarrow 0, ~ (n \rightarrow \infty )
\end{align*}
となるため，極限$B S$を${\mathcal B}(X)$内にもつ．
一方で，
\begin{align*}
  B S_n &= \left( I - ( I - B ) \right) S_n = S_n - (I - B) S_n \\
  &= \sum_{i=0}^n ( I - B)^i - \sum_{i=1}^{n+1} ( I - B)^i  \\
  &= I - ( I - B)^{n + 1}
\end{align*}
となり，定理の仮定より$\| I - B \|_{{\mathcal B}(X)} < 1$を持つため
\begin{align*}
  \| B S_n - I \|_{{\mathcal B}(X)} = \| ( I - B)^{n + 1} \|_{{\mathcal B}(X)} \le \| I - B \|_{{\mathcal B}(X)}^{n + 1} \rightarrow 0, ~ (n \rightarrow \infty )
\end{align*}
となるため，点列$(B S_n)$は極限$I$も${\mathcal B}(X)$内に持つ．
よって，極限の一意性より
\begin{align*}
  B S = I
\end{align*}
を得る．
${\mathcal R} ( I ) = X$であるため，${\mathcal R}(B S)=X$である．
その上，$X = {\mathcal R}(B S) \subset {\mathcal R}(B)$と${\mathcal R}(B) \subset X$となるため，${\mathcal R}(B) = X$となる．
さらに，$S \in {\mathcal B}(X)$であることから，
\begin{align*}
  {\mathcal D}(S) = {\mathcal R}(B) = X
\end{align*}
となる．

同様の議論を$S_n B \in {\mathcal B}(X)$について行うと
\begin{align*}
  S B = I
\end{align*}
と
\begin{align*}
{\mathcal D}(B) = {\mathcal R}(S) = X
\end{align*}
が得られる．
そのため，$B$は逆作用素を持ち，逆作用素$B^{-1} = S \in {\mathcal B}(X)$である．

また，
\begin{align*}
S_n = I + ( I - B) + ( I - B )^2 + \cdots + ( I - B )^n \rightarrow B^{-1}, ~ (n \rightarrow \infty)
\end{align*}
より
\begin{align*}
B^{-1} = I + (I - B) + (I - B)^2 + \cdots = \sum_{i = 0}^\infty (I - B)^i
\end{align*}
となる．

最後に
\begin{align*}
\left\| B^{-1} \right\|_{{\mathcal B}(X)} = \left\|  \sum_{i = 0}^\infty (I - B)^i \right\|_{{\mathcal B}(X)} \le \sum_{i=0}^\infty \| I - B \|^i
\end{align*}
となり，初項$1$，公比$\| I - B \|_{{\mathcal B}(X)} < 1$の総和より
\begin{align*}
\left\| B^{-1} \right\|_{{\mathcal B}(X)} \le \frac{1}{ 1 - \| I - B \|_{{\mathcal B}(X)} }
\end{align*}

\qed
\end{Proof}

Neumannの定理\ref{thm:nuemann_ser}は$B$をそのまま使う場合もありますが，$B = RA$のように二つの線形作用素の合成作用素として利用することも良くあります．
そのとき，$RA$が全単射であった場合における$R$と$A$の単射と全射に関する議論を行います．
そのため，次の定理もNeumannの定理\ref{thm:nuemann_ser}とセットで覚えておきましょう．

\begin{Thm}\label{thm:injective_surjective}
$X$と$Y$をBanach空間とする．
$A \in {\mathcal B}(X, Y)$, $R \in {\mathcal B}(Y, X)$とする．
もし$RA$が全単射ならば，$A$は単射であり，$R$は全射である．
\end{Thm}

\begin{Proof}
「$A$が単射」の証明\\
定理\ref{thm:linear_operator_injective11} (線形作用素に対する単射性(1))のii)を用いて証明する．
$u \in X$とし，$RA$が単射であることに注意すると
\begin{align*}
A u = 0 \Rightarrow RAu = 0 \Rightarrow u = 0
\end{align*}
よって，$A$は単射である．

「$R$が全射」の証明\\
$RA$が全射であるため任意の$g \in X$に対して，
\begin{align*}
RA u = g
\end{align*}
となる$u \in X$が存在する．
その上，$v = Au$とすると任意の$g \in X$に対して
\begin{align*}
R v = g
\end{align*}
となる$v \in Y$が存在するため，$R$は全射である．

\qed
\end{Proof}

Neumannの定理\ref{thm:nuemann_ser}と定理\ref{thm:injective_surjective}，及び，線形代数の次元定理を利用することで定理\ref{thm:basic_linear_problem}を証明することができます．
実際に$\| I - RA \| < 1$よりNeumannの定理\ref{thm:nuemann_ser}から$RA \in {\mathbb R}^n$は全単射になります．
そのうえで，定理\ref{thm:injective_surjective}から$R \in {\mathbb R}^n$は全射となり，$A \in {\mathbb R}^n$は単射になります．
さらに，線形代数の次元定理を利用することで，$R$と$A$がそれぞれ全単射となることもいえます．
よって，
\begin{align*}
	A x^* = b \Leftrightarrow& A (x^* - \hat{x}) = b - A \hat{x} \\
	 \Leftrightarrow& RA (x^* - \hat{x}) = R(b - A \hat{x}) \\
	 \Leftrightarrow& x^* - \hat{x}  = (RA)^{-1} R(b - A \hat{x})
\end{align*}
となるため
\begin{align*}
	\| x^* - \hat{x} \|_{\infty}  = \| (RA)^{-1} R(b - A \hat{x}) \|_{\infty} \le \| (RA)^{-1} \|_{\infty} \| R(b - A \hat{x}) \|_{\infty} \le \frac{ \| R(b - A \hat{x}) \|_{\infty} }{1 - \| I - RA \|_{\infty}} 
\end{align*}
が成立します．

%
%----------------------------------------------------------------------------------------------------------------------------------------------
%----------------------------------------------------------------------------------------------------------------------------------------------
%----------------------------------------------------------------------------------------------------------------------------------------------
\newpage
\section{非線形解析の基礎}\label{sec:nonlinear_analysis}
本書では線形，非線形，有限次元，無限次元を問わず方程式のコンピュータで得られた解の品質を保証するための定理を紹介することが目的です．
そのためには，無限次元空間上の非線形作用素を解析するための道具が必要になります．
有限次元の言葉を借りて説明すると，全微分，Riemann積分，区間上のベクトル値関数に対する微分積分学の基本定理(Fundamental theorem of calculus)が必要になります．
Banach空間上においてはFr\'echet微分，Bochner積分，区間上のBanach空間値関数に対する微分積分学の基本定理が必要になります．


%----------------------------------------------------------------------------------------------------------------------------------------------
\subsection{Fr\'echet微分}\label{sec:define_Frechet_diff}
本節ではBanach空間上の作用素に対する微分としてFr\'echet微分を紹介します．
ベクトル空間${\mathbb R}^n$から${\mathbb R}^m$への写像に対する微分は方向微分と全微分がありましたね．
Fr\'echet微分は全微分の拡張にあたります．
本書では取り扱いませんが，方向微分の拡張はG\^ateaux微分と呼ばれます．


\begin{defi}[Fr\'echet微分]\label{def:frechet_diff}
$X, Y$をBanach空間とし，開部分集合$U \subset X$とする．
定義域を${\mathcal D}(f) = U$とする$U$から$Y$への作用素$f$は$U$上で連続とする．
ある点$v \in U$に対し，$v + h \in U$となる任意の$h \in X$について
\begin{align*}
\frac{ \| f(v + h) - f(v) - f'[v] h \|_{Y} }{\| h \|_{X}} \rightarrow 0, ~ (h \rightarrow 0)
\end{align*}
を満たす線形作用素$f'[v] \in {\mathcal B}(X, Y)$が存在するとき，作用素$f$は点$v$においてFr\'echet微分可能といい，$f'[v] \in {\mathcal B}(X, Y)$を$f$の点$v$におけるFr\'echet微分と呼ぶ.
\end{defi}

Fr\'echet微分の定義\ref{def:frechet_diff}はランダウの記号$o$を用いると
\begin{align*}
  f(v + h) - f(v) =  f'[v] h + o( h )
\end{align*}
となる$o(h)$が
\begin{align*}
  \frac{\| o( h ) \|_{Y}  }{ \| h \|_{X} } \rightarrow 0, ~ (h \rightarrow 0)
\end{align*}
を満たすときです．

開部分集合$U$の任意の元でFr\'echet微分可能なとき，$f'$を$U$から${\mathcal B}(X, Y)$への作用素，すなわち，$f': U \rightarrow {\mathcal B}(X, Y)$と見ることもできます．
実際に$v \in U$について$f'[v] \in {\mathcal B}(X, Y)$となることは定義から明らかです．

Fr\'echet微分の定義\ref{def:frechet_diff}では，$U$が開集合であることが重要な点の一つです．
ノルム空間の開集合の定義\ref{def:App:OpenSet}を思い出すと，$U$の任意の元$v$に対して開球$B_X(v, r) \subset U$となる$r>0$が必ず存在します．
すなわち，開部分集合$U$の任意の元$v$に対して，ありとあらゆる方向にほんの少し摂動させても，まだ$U$に属することができることを意味します．
よって，「$v + h \in U$となる任意の$h \in X$」でいう$h$は$U$におさまる範囲で$v$をありとあらゆる方向に摂動させることを意味します．
実際に${\mathbb R}$で考えてみても$v \in (0,1)$であれば，どんな$v$でも正の方向に摂動しても，負の方向に摂動しても$(0,1)$に属する元は存在します．
しかし，$v \in [0,1)$としてしまうと，$v = 0$のとき負の方向に摂動すると$[0,1)$に属する元は存在しません．


また，\ref{chap:linear_eq}章や\ref{chap:eigenvalue}章で扱ったように計算の品質保証では閉区間あるいは閉球の中に解が存在することを証明します．
その際に，解が存在するであろう範囲の候補となる閉部分集合上でFr\'echet微分が可能であることを仮定します．
これは，「解が存在するであろう範囲の候補となる閉部分集合」よりも広い開部分集合でFr\'echet微分が可能なことを意味します．
感覚的には$f(x) = \sqrt{x}$を考えた際に$X = {\mathbb R}$, $U = (0, \infty)$として，$U$上で微分可能であるが，あえて「$f$は閉区間$[1, 2]$上で微分可能である」と言うことと同じです．

続いて，Fr\'echet微分の性質を紹介します．

\begin{Thm}[Fr\'echet微分の線形性]\label{thm:frechet_linear_arith}
$X, Y$をBanach空間とし，開部分集合$U \subset X$とする．
定義域を${\mathcal D}(f_1) = {\mathcal D}(f_2)  = U$とする$U$から$Y$への作用素$f_1$と$f_2$が，各々$v \in U$でFr\'echet微分可能とし，
各々のFr\'echet微分を$f_1'[v]$と$f_2'[v]$と表記する．
${\mathcal D}(g)  = U$とする$U$から$Y$への作用素$g$を
\begin{align*}
g(u) = \alpha f_1(u) + \beta f_2(u), ~ u \in U, ~ \alpha, \beta \in {\mathbb K}
\end{align*}
とすると，作用素$g$も$v \in U$にてFr\'echet微分可能である．
その上，$g$の$v \in U$におけるFr\'echet微分を$g'[v]$と表記すると
\begin{align*}
g'[v] = \alpha f_1'[v] + \beta f_2'[v]
\end{align*}
が成立する．
\end{Thm}

\begin{Proof}
作用素$f_1$と$f_2$が$v$においてFr\'echet微分可能であるため，ランダウの記号を用いた表記
\begin{align*}
  \frac{\| o_1( h ) \|_{Y}  }{ \| h \|_{X} } \rightarrow 0, ~ (h \rightarrow 0)
\end{align*}
となる
\begin{align*}
  f_1(v + h) - f_1(v) =  f_1'[v] h + o_1( h )
\end{align*}
と，
\begin{align*}
  \frac{\| o_2( h ) \|_{Y}  }{ \| h \|_{X} } \rightarrow 0, ~ (h \rightarrow 0)
\end{align*}
となる
\begin{align*}
  f_2(v + h) - f_2(v) =  f_2'[v] h + o_2( h )
\end{align*}
を持つ．

\begin{align*}
g(v + h) &= \alpha f_1(v + h) + \beta f_2( v + h ) \\
&= \alpha \left( f_1(v) + f_1'[v] h + o_{f_1}( h ) \right) + \beta \left(  f_2(v) + f_2'[v] h + o_{f_2}( h ) \right) \\
&= \left( \alpha f_1(v) + \beta f_2(v) \right) + \left( \alpha f_1'[v] + \beta f_2'[v] \right) h +  \left( \alpha o_{f_1}( h ) + \beta  o_{f_2}( h ) \right) \\
&= g(v) +  \left( \alpha f_1'[v] + \beta f_2'[v] \right) h +  \left( \alpha o_{f_1}( h ) + \beta  o_{f_2}( h ) \right)
\end{align*}
となる．
その上，
\begin{align*}
\frac{ \| \alpha o_{f_1}( h ) + \beta  o_{f_2}( h ) \|_{Y} }{ \| h \|_X } \le \frac{ | \alpha | \| o_{f_1}( h ) \|_{Y} }{\| h \|_X} + \frac{ | \beta |  \|  o_{f_2}( h ) \|_{Y} }{ \| h \|_X } \rightarrow 0, ~ (h \rightarrow 0)
\end{align*}
となるため，作用素$g$は$v \in U$においてFr\'echet微分可能で，$g$のFr\'echet微分$g'[v]$は
\begin{align*}
g'[v] = \alpha f_1'[v] + \beta f_2'[v]
\end{align*}
となる．

\qed
\end{Proof}



\begin{Thm}[Fr\'echet微分の連鎖律]\label{thm:frechet_chain}
$X, Y, Z$をBanach空間とする．
開部分集合を$U_1 \subset X$とし，定義域を${\mathcal D}(f_1) = U_1$とする$U_1$から$Y$への作用素$f_1$が，$v \in U_1$でFr\'echet微分可能とする．
また，$U_2 \subset {\mathcal R}(f_1) \subset Y$とし，定義域を${\mathcal D}(f_2) = U_2$とする$U_2$から$Z$への作用素$f_2$が，$f_1(v) \in U_2 \subset {\mathcal R}(f_1)$でFr\'echet微分可能とする．
各々のFr\'echet微分を$f_1'[v]$と$f_2'[f_1(v)]$と表記する．
${\mathcal D}(g)  = U_1$とする$U_1$から$Z$への作用素$g$を
\begin{align*}
g(u) = f_2( f_1(u) ), ~ u \in U_1
\end{align*}
とすると，作用素$g$も$v \in U_1$にてFr\'echet微分可能である．
その上，$g$の$v \in U_1$におけるFr\'echet微分を$g'[v]$と表記すると
\begin{align*}
g'[v] = f_2'[ f_1( v ) ] f_1'[ v ]
\end{align*}
が成立する．
\end{Thm}

\begin{Proof}
作用素$f_1$が$U_1$でFr\'echet微分可能であるため，ランダウの記号を用いた表記
\begin{align*}
  \frac{\| o_1( h_1 ) \|_{Y}  }{ \| h_1 \|_{X} } \rightarrow 0, ~ (h_1 \rightarrow 0)
\end{align*}
となる
\begin{align*}
  f_1(v + h_1) - f_1(v) =  f_1'[v] h_1 + o_1( h_1 )
\end{align*}
となる．
同様に作用素$f_2$が$U_2$でFr\'echet微分可能であるため
\begin{align*}
  \frac{\| o_2( h_2 ) \|_{Z}  }{ \| h_2 \|_{Y} } \rightarrow 0, ~ (h_2 \rightarrow 0)
\end{align*}
となる
\begin{align*}
  f_2( f_1(v) + h_2 ) - f_2( f_1(v) ) =  f_2'[ f_1(v) ] h_2 + o_2( h_2 )
\end{align*}
をもつ．

\begin{align*}
  f_2( f_1( v + h_1 ) ) = f_2( f_1(v) + f_1'[v] h_1 + o_1( h_1 ) )
\end{align*}
となる．
$h_2 =  f_1'[v] h_1 + o_1( h_1 )$とおくと
\begin{align*}
  f_2( f_1( v + h_1 ) ) &= f_2( f_1(v) + f_1'[v] h_1 + o_1( h_1 ) ) \\
&= f_2( f_1(v) + h_2) \\
&=  f_2( f_1(v) ) + f_2'[ f_1(v) ] h_2 + o_2( h_2 ) \\
&=  f_2( f_1(v) ) + f_2'[ f_1(v) ] ( f_1'[v] h_1 + o_1( h_1 ) ) + o_2( h_2 ) \\
&=  f_2( f_1(v) ) + f_2'[ f_1(v) ] f_1'[v] h_1 + \left( f_2'[ f_1(v) ] o_1( h_1 ) + o_2( h_2 ) \right)
\end{align*}
となる．
その上，Fr\'echet微分の定義より$f_2'[ f_1(v) ] \in {\mathcal B}(Y, Z)$であるため，$\| f_2'[ f_1(v) ] \|_{{\mathcal B}(Y, Z)} \le M$となる定数$M > 0$が存在する．
また，
\begin{align*}
\frac{ \| f_2'[ f_1(v) ] o_1( h_1 ) + o_2( h_2 ) \|_{ Z } }{ \| h_1 \|_{ X } } &\le \frac{ \| f_2'[ f_1(v) ] o_1( h_1 ) \|_{Z} }{ \| h_1 \|_{ X } } + \frac{ \| o_2( h_2 ) \|_{ Z } }{ \| h_1 \|_{X} } \\
&\le M \frac{ \| o_1( h_1 ) \|_{Y} }{ \| h_1 \|_{ X } } + \frac{ \| o_2( h_2 ) \|_{ Z } }{ \| h_2 \|_{Y} } \frac{ \| h_2 \|_{Y} }{ \| h_1 \|_{X} }
\end{align*}
となり，
\begin{align*}
\frac{ \| h_2 \|_{Y} }{ \| h_1 \|_{X} } &= \frac{ \|  f_1'[v] h_1 + o_1( h_1 ) \|_{Y} }{ \| h_1 \|_{X} } \\
&\le M + \frac{ \| o_1( h_1 ) \|_{Y} }{ \| h_1 \|_{X} } 
\end{align*}
から
\begin{align*}
& \frac{ \| f_2'[ f_1(v) ] o_1( h_1 ) + o_2( h_2 ) \|_{ Z } }{ \| h_1 \|_{ X } } \\
&\le M \frac{ \| o_1( h_1 ) \|_{Y} }{ \| h_1 \|_{ X } } + \frac{ \| o_2( h_2 ) \|_{ Z } }{ \| h_2 \|_{Y} } \left( M + \frac{ \| o_1( h_1 ) \|_{Y} }{ \| h_1 \|_{X} }  \right) \rightarrow 0, ~ h_1, h_2 \rightarrow 0
\end{align*}
となるため，作用素$g$は$v \in U_1$においてFr\'echet微分可能で，$g$のFr\'echet微分$g'[v]$は
\begin{align*}
g'[v] = f_2'[ f_1( v ) ] f_1'[ v ]
\end{align*}
となる．

\qed
\end{Proof}




%------------------------------------------------------------------------------------------------------------------------
\subsection{Bochner積分}\label{sec:bochner_integral}
本節ではBochner積分を紹介します．
本書では，定理には証明をつけるスタンスですが，本節のみ証明を省きます．
理由としてはLebesgue積分のBanach空間に値をとる関数への拡張であり，測度論からLebesgue積分論，Banach空間に値をとる関数の可測性を経てBochner積分が定義され，定理が証明されるので，本書のゆったりペースでは1冊かけて説明しなくてはならなくなってしまいます．
測度論及びLebesgue積分論を知っているという前提で，証明や詳細を知りたい場合は，宮寺功「関数解析」の第6章を参照してください．

本節ではBochner積分のエッセンスのみ提示します．
$f$を実数の閉区間$[a, b]$から${\mathbb R}$の写像としたとき，$f$の$[a, b]$上のRiemann積分は
\begin{align*}
\int_a^b f(x) dx
\end{align*}
のように微積分で習っていると思います．
Bochner積分では，\\
\noindent ~~ 1. 区間$[a, b]$上ではなく一般的な測度空間$S$上の積分を考える \\
\noindent ~~ 2. $f$をBanach空間$X$への作用素にする\\
ことです．
すなわち，$f$を測度空間$S$からBanach空間$X$への作用素としたときの，$f$の$S$上の積分を定義することがBochner積分の趣旨です．


しかし，本書では，具体的には一般的な測度空間$S$を扱う必要がなく，実数の閉区間$[a, b]$からBanach空間$X$への作用素$f$の積分のみで十分です．
その上，結果的に言うと，いわゆる，積分の線形性や交換則などのRiemann積分で習った性質もちゃんと受け継がれています．
また，区間$[a, b]$からベクトル${\mathbb R}^n$への写像を区間上のベクトル値関数と呼ぶ風習にならって，区間$[a, b]$からBanach空間$X$への作用素を区間上のBanach空間値関数と呼びます．


集合$S$のいくつかの部分集合からなる集合族を$\Sigma$を完全加法集合体とし，$\mu$を$\Sigma$上で定義される測度とします．
測度と完全加法集合体について細かい定義は省略しますが，測度とは面積や体積の一般化された概念であり，完全加法集合体とは測度を定義するには十分な性質を備えた集合の集まりです．
集合，集合族，測度の三つの組$(S, \Sigma, \mu)$を測度空間と呼びます．

では，まず，単純関数を定義します．
単純関数のイメージは区間$[a, b]$から${\mathbb R}$への関数が，区間を$[a, c_1), [c_1, c_2), \cdots [c_n, b]$のように分割したときに，各々の分割した区間上で定数になる場合の一般化になります．

\begin{defi}[単純関数]\label{def:simple_function}
$(S, \Sigma, \mu)$を測度空間とし，可測集合$A_i \in \Sigma, ~ i = 1, 2, \cdots, n$を
\begin{align*}
A_i \cap A_j, ~ i \neq j ~~  ~ \mbox{が空集合}
\end{align*}
かつ
\begin{align*}
S = \bigcup_{i = 1}^{ \infty } A_i
\end{align*}
とする．
$X$をBanach空間とし，定義域が$S$である$S$から$X$への作用素$f$が各$A_i$上でBanach空間$X$としての定数をとるとき，$f$を単純関数と呼ぶ．
\end{defi}

続いて，単純関数のBochner積分を定義します．

\begin{defi}[単純関数のBochner積分]\label{def:bochner_integral_for_simple_function}
$(S, \Sigma, \mu)$を測度空間とし，$X$をBanach空間とする．
$f$を$S$から$X$への単純関数とする．
すなわち，
\begin{align*}
f(s) = f_n, ~ (s \in A_n) 
\end{align*}
となるBanach空間$X$の定数$f_n \in X$と
\begin{align*}
A_i \cap A_j, ~ i \neq j ~~  ~ \mbox{が空集合}
\end{align*}
かつ
\begin{align*}
S = \bigcup_{i = 1}^{ \infty } A_i
\end{align*}
となる可測集合$A_i \in \Sigma, ~ i = 1, 2, \cdots, n$が存在する．
そのとき，
\begin{align*}
 \| f(s) \|_{X} \mbox{が}S\mbox{上でLebesgue可積分}
\end{align*}
のとき，$f(s)$は$S$上でBochner可積分であるといい，$f(s)$の$S$上のBochner積分を
\begin{align*}
 \int_S f(s) d \mu := \sum_{i = 1}^{\infty} f_n \mu(A_i) \in X
\end{align*}
と定義する．
\end{defi}

$\mu(A_i)$は集合$A_i$の測度です．
厳密ではない言葉で表現すると，集合$A_i$の面積に相当するものですので，$\mu(A_i)$の定義には$\mu(A_i) \ge 0$が含まれます．
そのため，単純関数のBochner積分の定義より
\begin{align*}
 \left\| \int_S f(s) d \mu \right\|_{X} \le \int_S \| f(s) \|_{X} d \mu
\end{align*}
を満たします．
(上式の右辺になればLebesgue積分で十分ですね．)

つづいては一般的な関数のBochner積分を定義します:

\begin{defi}[Bochner積分]\label{def:bochner_integral}
$(S, \Sigma, \mu)$を測度空間とし，$X$をBanach空間とする．
定義域${\mathcal D}(f) = S$となる$S$から$X$への作用素を$f$とする．
$s \in S$に対してほとんど至るところで
\begin{align*}
f_n(s) \rightarrow f(s), ~  (n \rightarrow \infty) 
\end{align*}
かつ
\begin{align*}
\int_S \| f(s) - f_n(s) \|_{X} d \mu \rightarrow 0, ~ ( n \rightarrow \infty)
\end{align*}
となるような$S$上でBochner可積分な単純関数の列$\{ f_n(s) \}$が存在するとき，$f(s)$は$S$上でBochner可積分であるという．
その上，$f(s)$の$S$上のBochner積分を$\int_S f_n(s) d$の極限，すなわち
\begin{align*}
\int_S f_n(s) d \mu \rightarrow \int_S f(s) d \mu , ~ ( n \rightarrow \infty)
\end{align*}
として定義する．
\end{defi}

これでBochner積分が定義できました．
続いて，Bochner積分の線形性や有界作用素との関係など本書で使う性質を証明なしで紹介します．

\begin{Thm}[Bochner積分のノルム評価]
	$(S, \Sigma, \mu)$を測度空間とし，$X$をBanach空間とする．
	$f$を$S$から$X$への$S$上のBochner可積分な関数とする．
	そのとき
	\begin{align*}
		\left\| \int_S f(s) d \mu \right\|_{X} \le \int_S \| f(s) \|_{X} d \mu
	\end{align*}
	となる．
\end{Thm}

\begin{Thm}[Bochner積分の線形性]
	$(S, \Sigma, \mu)$を測度空間とし，$X$をBanach空間とする．
	$f$と$g$を$S$から$X$への$S$上のBochner可積分な関数とする．
	$\alpha, \beta \in {\mathbb K}$としたとき，
	\begin{align*}
		\int_S \left( \alpha f(s) + \beta g(s) \right) d \mu = \alpha \int_S f(s) d \mu + \beta \int_S g(s) d \mu
	\end{align*}
	となる．
\end{Thm}

\begin{Thm}[Bochner積分の線形性]
	$(S, \Sigma, \mu)$を測度空間とし，$X$をBanach空間とする．
	$f$と$g$を$S$から$X$への$S$上のBochner可積分な関数とする．
	$\alpha, \beta \in {\mathbb K}$としたとき，
	\begin{align*}
		\int_S \left( \alpha f(s) + \beta g(s) \right) d \mu = \alpha \int_S f(s) d \mu + \beta \int_S g(s) d \mu
	\end{align*}
	となる．
\end{Thm}

\begin{Thm}[有界線形作用素とBochner積分の関係]
	$(S, \Sigma, \mu)$を測度空間とし，$X$と$Y$をBanach空間とする．
	$f$を$S$から$X$への$S$上のBochner可積分な関数とする．
	$B \in {\mathcal B}(X, Y)$とすると
	$\alpha, \beta \in {\mathbb K}$としたとき，
	\begin{align*}
		\int_S B f(s) d \mu = B \int_S f(s) d \mu
	\end{align*}
	となる．
\end{Thm}

続いて，測度空間$(S, \Sigma, \mu)$の$S$を閉区間$[a, b] \subset {\mathbb R}$のように限定した場合の性質を紹介します:

\begin{defi}[閉区間上のBochner積分]\label{def:bochner_integral_for_interval}
$X$をBanach空間とし，$f$を閉区間$[a, b]$から$X$への$[a,b]$上のBochner可積分な関数とする．
そのとき，$[a,b]$上のBochner積分を
	\begin{align*}
		\int_a^b f(s) d s := \int_{[a,b]} f(s) d \mu
	\end{align*}
と記述する．
\end{defi}

\begin{Thm}[閉区間上の連続作用素とBochner積分]\label{thm:Bochnerable}
	$X$をBanach空間とし，$f$を閉区間$[a, b]$から$X$への$[a,b]$上のBochner可積分な関数とする．
	$f$が任意の$t \in [a, b]$において連続ならば$[a, b]$上でBochner可積分である．
\end{Thm}

%------------------------------------------------------------------------------------------------------------------------
\subsection{閉区間上のBanach空間値関数に対する微分積分学の基本定理}\label{sec:Fund_theorem}
本章の最後にFr\'echet微分と閉区間上のBochner積分の関係性を示す微分積分学の基本定理を導出します．


\begin{Thm}[Fr\'echet微分と閉区間上におけるBochner積分に対する微分積分学の基本定理(i)]\label{the:fund_theorem_i}
$X$をBanach空間とする．
$f$を閉区間$[a, b]$から$X$への$[a,b]$上で連続な関数とする．
$t \in (a, b )$とし，関数$g$を
\begin{align*}
	g( t ) := \int_a^t f(s) ds
\end{align*}
とすると，$g$は$t \in (a, b )$においてFr\'echet微分可能である．
その上，$g$の$t$のFr\'echet微分を$g'[t]$とすると
\begin{align*}
	g'[t] = f(t)
\end{align*}
となる．
\end{Thm}

\begin{Proof}
$f$は任意の$t \in [a, b]$において連続であるため，
\begin{align*}
	\forall \varepsilon > 0 ~ \exists \delta > 0 ~ | s - t | < \delta ~ \mbox{となる} ~ \forall s \in [a, b] ~ \mbox{に対して} ~ \| f(s) - f(t) \|_{X} < \varepsilon
\end{align*}
が成立する．
その上，$g$の$t$のFr\'echet微分可能か定義にあてはめると
\begin{align*}
	\frac{ \left\| g(t + h) - g(t) - f(t)h \right\|_{X} }{ |h| } &= \frac{ \left\| \int_{a}^{t+h} f(s) ds - \int_{a}^{t} f(s) ds - f(t)h \right\|_{X} }{ |h| } \\
	&= \frac{ \left\| \int_{t}^{t+h} f(s) ds - f(t)h \right\|_{X} }{ |h| } \\
	&= \frac{ \left\| \int_{t}^{t+h} f(s) ds - \left( f(t)( t+h) - f(t)t \right) \right\|_{X} }{ |h| } \\
	&= \frac{ \left\| \int_{t}^{t+h} f(s) ds - \int_{t}^{t+h} f(t) ds \right\|_{X} }{ |h| } \\
	&= \frac{ \left\| \int_{t}^{t+h} \left( f(s)  - f(t) \right) ds \right\|_{X} }{ |h| } \\
	&\le \frac{ \int_{t}^{t+h} \left\| f(s)  - f(t) \right\|_{X} ds }{ |h| } \\
	&< \frac{ \varepsilon   \int_{t}^{t+h}  ds }{ |h| } = \varepsilon 
\end{align*}
となる．
その上，$\varepsilon > 0$は任意であるため，
\begin{align*}
	\frac{ \left\| g(t + h) - g(t) - f(t)h \right\|_{X} }{ |h| } \rightarrow 0, ~ (h \rightarrow 0)
\end{align*}
となるため，$g$はFr\'echet微分可能で$g$の$t$におけるFr\'echet微分は$f(t)$となる．

\qed
\end{Proof}


\begin{Thm}[Fr\'echet微分と閉区間上におけるBochner積分に対する微分積分学の基本定理(ii)]\label{thm:fund_theorem_ii}
$X$をBanach空間とする．
$f$を閉区間$[a, b]$から$X$への$[a,b]$上で連続な関数とし，その上，任意の$s \in (a, b) $においてFr\'echet微分可能とする．
$f$の$s \in (a, b)$におけるFr\'echet微分を$f'[s]$とし，$f': [a, b] \rightarrow {\mathcal B}({\mathbb R},X)$で連続
\footnote{ 
$f'[v] \in {\mathcal B}(X, Y)$が$h \in Y$で連続とは
\[
	\forall \varepsilon > 0, ~ \exists > 0, ~ \| h_n - h \|_X < \delta ~ \mbox{となる} ~ \forall h_n \in X ~ \mbox{に対して} ~ \| f'[v] h_n - f'[v] h \|_{Y} < \varepsilon
\]
が成立することです．
その一方で，$f' : [a, b] \rightarrow {\mathcal B}(X, Y)$が$h \in [a,b]$で連続とは
\[
	\forall \varepsilon > 0, ~ \exists > 0, ~ | h_n - h | < \delta ~ \mbox{となる} ~ \forall h_n \in {\mathbb R} ~ \mbox{に対して} ~ \| f'[h_n] - f'[h]  \|_{{\mathcal B}(X,Y)} < \varepsilon
\]
が成立することです．
}
であるとすると，
\begin{align*}
	\int_{a}^{b} f'[s] ds = f(b) - f(a)
\end{align*}
となる．
\end{Thm}

\begin{Proof}
	$f'$が$[a, b]$から${\mathcal B}({\mathbb R},X)$で連続であるため定理\ref{thm:Bochnerable}から$f'[s], ~ s \in [a, b]$でBochner可積分である．
	その上，$g$を
	\begin{align*}
		g(t) = \int_{a}^{t} f'[s] ds
	\end{align*}
	とすると定理\ref{the:fund_theorem_i}から，$g$は$t \in (a, b)$においてFr\'echet微分可能で，
	\begin{align*}
		g'[t] = f'[t]  ~ \forall t \in (a, b)
	\end{align*}	
	となる．
	\begin{align*}
		F(t) = g(t) - f(t)
	\end{align*}	
	とおくと，$g$も$f$も$t \in (a, b)$においてFr\'echet微分可能であることから，
	\begin{align*}
		F'[t] = g'[t] - f'[t] = 0, ~ \forall t \in (a, b)
	\end{align*}	
	となる．
	また，$F'[t]$がゼロであることから，Fr\'echet微分の定義により$F(t)$は$t$に依存しない作用素になる．
	よって
	\begin{align*}
		F(t) = g(t) - f(t) = c
	\end{align*}	
	となる$t \in [a, b]$によって変化しない$c \in X$が存在する．
	その上，上の式の$t$に$a$を代入すると
	\begin{align*}
		g(a) - f(a) = 0 - f(a) = c
	\end{align*}	
	となるため，$c = -f(a)$である．
	よって，$g(t) - f(t) = c$に注意すると
	\begin{align*}
		\int_{a}^{t} f'[s] ds = g(t) = f(t) + c = f(t) - f(a)
	\end{align*}	
	が成り立つ．
	ゆえに$t = b$とすると
	\begin{align*}
		\int_{a}^{b} f'[s] ds = f(b) - f(a)
	\end{align*}	
	となる．

\qed
\end{Proof}

\begin{Thm}[Fr\'echet微分と閉区間上のBochner積分に対する微分積分学の基本定理(iii)]\label{thm:fundamental_theorem_of_calculus}
$X, Y$をBanach空間とし，$U$を$X$の開部分集合とする．
作用素$F:U \rightarrow Y$とし，任意の$v \in U$においてFr\'echet微分可能であるとする．
$W$を$W \subset U$となる(開，閉を問わない)部分集合とする．
任意の$u, v \in W$について
\begin{align*}
	w(t) = (1 - t) u + t v, ~ t \in [0,1]
\end{align*}
が$W$に属すると仮定する．
$F': U \rightarrow {\mathcal B}(X,Y)$が$W$上で連続であると仮定する．
そのとき，
\begin{align*}
	F(v)  - F(u) = \int_0^1 F'[(1 - t) u + t v](v - u)dt
\end{align*}
となる．
\end{Thm}

\begin{Proof}
$\tilde{w}: (1 - t) u + t v, ~ t \in {\mathbb R}$となる$\tilde{w}: {\mathbb R} \rightarrow X$としたとき，$\tilde{w}$は${\mathbb R}$上でFr\'echet微分可能で
\[
	\tilde{w}'[t] = v - u
\]
となる．
そのうえで，$\tilde{w}' : {\mathbb R} \rightarrow {\mathcal B}({\mathbb R}, X)$は$\| \tilde{w}'[s] - \tilde{w}'[t] \|_{X} = 0$となるため，${\mathbb R}$上で連続である．
よって，$\tilde{w}$の定義域を$[0, 1]$に制限した$w:[0, 1] \rightarrow W$は$[0, 1]$上でFr\'echet微分可能であり，$\tilde{w}'$の定義域を$[0,1]$にしたFr\'echet微分$w' : [0, 1] \rightarrow {\mathcal B}({\mathbb R}, X)$は$[0, 1]$上で連続である．

$f: [0, 1] \rightarrow Y$を
\begin{align*}
	f(t) := F(w(t))
\end{align*}
とする．
$W \subset U$から作用素$F:U \rightarrow Y$が任意の$v \in W$においてFr\'echet微分可能であり，$f$のFr\'echet微分は連鎖律定理\ref{thm:frechet_chain}より
\begin{align*}
	f'[t] = F'[w(t)]w'[t] = F'[w(t)]( v - u)
\end{align*}
となるため，任意の$[0, 1]$上でFr\'echet微分可能である．
また，仮定から$F': U \rightarrow {\mathcal B}(X, Y)$が$W$上で連続していることより，$f':[0, 1] \rightarrow {\mathcal B}({\mathbb R}, Y)$は$[0, 1]$上で連続となる．
よって，$f$に対して定理\ref{thm:fund_theorem_ii}を利用でき，
\begin{align*}
	f(1) - f(0) = \int_0^1 f'[t] dt
\end{align*}
となる．
ゆえに，
\begin{align*}
	F(v)  - F(u) = f(1) - f(0) = \int_0^1 f'[t] dt = \int_0^1 F'[w(t)]( v - u) dt
\end{align*}
となり，題意は示せた．

\qed
\end{Proof}


%
%----------------------------------------------------------------------------------------------------------------------------------------------
%----------------------------------------------------------------------------------------------------------------------------------------------
%----------------------------------------------------------------------------------------------------------------------------------------------
\newpage
\section{方程式の解の品質保証の準備}
本章では，数値計算を用いて得られる方程式の解を品質保証するための基本的な定理を紹介します．
品質保証するための基本定理はBanach空間の抽象的な関数方程式における解の存在定理が根幹にあり，不動点定理やNewton-Kantorovichの定理を品質保証のために使いやすくした定理だと思ってください．
今まで導出してきたNeumann級数の定理\ref{thm:nuemann_ser}，Fr\'echet微分(定義\ref{def:frechet_diff})，Fr\'echet微分と閉区間上のBochner積分に対する微分積分学の基本定理(iii)(定理\ref{thm:fundamental_theorem_of_calculus})のほかに，Banachの不動点定理を利用します．
そのため，まず，Banachの不動点定理の紹介から始めます．



%----------------------------------------------------------------------------------------------------------------------------------------------
\subsection{Banachの不動点定理}
\begin{defi}[不動点]
$X$を係数体が${\mathbb K}$のBanach空間とする．
$M$は空でない閉集合で$M \subset X$を満たすとする．
$A$を$M$から$M$への写像とする．
$x \in M$が$A$の不動点であるとは，$x$が
\begin{align*}
x = Ax
\end{align*}
を満たすことである．
\end{defi}
         
\begin{defi}[縮小写像]
$X$を係数体が${\mathbb K}$のBanach空間とする．
$M$は空でない閉集合で$M \subset X$を満たすとする．
写像$A: M \rightarrow M$が$k$次の縮小写像であるとは，$0 \le k < 1$を満たす定数$k$が存在し，
$\forall x,y \in M$について
\begin{align*}
\| A x - A y \| \le k \| x - y\|
\end{align*}
を満たすことである．
\end{defi}

\begin{Thm}[Banachの不動点定理]\label{thm:Banach}
$X$を係数体が${\mathbb K}$のBanach空間とする．
$M$は空でない閉集合で$M \subset X$を満たすとする．
$A$は$M$から$M$への$k$次の縮小写像とする．
そのとき，問題
\begin{equation}
\text{Find} ~ u \in M ~ \text{s.t.} ~ u = Au
\label{Eq:FixedPointEq}
\end{equation}
は真の解$u^*$を$M$内にただ一つ持つ．
即ち，写像$A$は$M$上にただ一つ不動点$u^*$を持つ．
\end{Thm}

\begin{Cor}\label{cor:Banach}
$X$,$M$,$A$,$k$は定理\ref{thm:Banach}で定義されている集合，写像及び定数とする．
$u_0$を閉集合$M$の元として与えられていると仮定する．
そのとき反復法
\begin{align*}
u_{n+1} = A u_n, ~ n = 0,1,\cdots
\end{align*}
によって得られる点列$(u_n)$は(\ref{Eq:FixedPointEq})を満たす真の唯一解$u^{*}$に収束する．
さらに，$\forall n \in {\mathbb N}$について不等式
\begin{align}
\| u_n - u^* \| &\le k^n (1-k)^{-1} \| u_1 - u_0 \| \label{Ine:Apri} \\
\| u_{n+1} - u^* \| &\le k (1-k)^{-1} \| u_{n+1} - u_n \| \label{Ine:Apos} \\
\| u_{n+1} - u^* \| &\le k \| u_n - u^* \| \label{Ine:rate}
\end{align}
が成立する．
\end{Cor}

\begin{Rem}
(\ref{Ine:Apri})，(\ref{Ine:Apos})，(\ref{Ine:rate})はそれぞれ事前誤差評価，事後誤差評価，収束率と呼ばれる．
\end{Rem}

%----------------------------------------------------------------------------------------------------------------------------------------------
\begin{Proof}[Banachの不動点定理]
$u_0$を閉集合$M$の元として与えられていると仮定する．
点列$(u_n)$は反復法
\begin{equation}
u_{n+1} = A u_n, ~ n = 0,1,\cdots \label{pr:Eq:ita}
\end{equation}
によって得られる．
そのとき，証明のプロセスは次のように考える:\\
I) $(u_n)$がCauchy列になること，さらにBanach空間の完備性を使うことで，$u_n \rightarrow u, ~ n \rightarrow \infty$となる$u$が$X$内に存在することを示す．\\
II) $u$が(\ref{Eq:FixedPointEq})を満たす真の解$u^*$と一致することを示す(解の存在性)．\\
III) 真の解$u^*$が$M$内で一意であることを示す．\\
\\
I)

(\ref{pr:Eq:ita})より
\begin{align*}
\| u_n - u_{n+1} \| = \| A u_{n-1} - A u_n \|
\end{align*}
となる．
仮定より$A$は$k$次の縮小写像であるため，
\begin{align*}
\| A u_{n-1} - A u_n \| \le k \| u_{n-1} - u_n \|
\end{align*}
となる定数$k$が存在する．
同様に$\| u_{n-1} - u_n \|$に(\ref{pr:Eq:ita})と$A$の縮小写像の性質を使うと
最終的に
\begin{equation}
\| u_n - u_{n+1} \| \le k^n \| u_0 - u_1 \|
\label{prof:1}
\end{equation}
を得る．


次に三角不等式(定義\ref{def:App:norm} (iv))より，$n=0,1,2,\cdots, ~ m > n $について
\begin{align}
\| u_n - u_m \| &= \| (u_n - u_{n+1}) + (u_{n+1} - u_{n+2} ) + \cdots + ( u_{m-1} - u_m) \| \nonumber \\
&\le \| u_n - u_{n+1} \| + \| u_{n+1} - u_{n+2} \| + \cdots + \| u_{m-1} - u_m \|
\label{prof:3}
\end{align}
となる．
上の式に(\ref{prof:1})を適用すると
\begin{align}
\| u_n - u_m \| &\le \| u_n - u_{n+1} \| + \| u_{n+1} - u_{n+2} \| + \cdots + \| u_{m-1} - u_m \| \nonumber \\
&\le k^n \| u_0 - u_1 \| + k^{n+1} \| u_0 - u_1 \| + \cdots + k^{m-1} \| u_0 - u_1 \| \nonumber \\
&= k^n (1+k+\cdots+k^{m-n-1}) \| u_0 - u_1 \|
\label{prof:2}
\end{align}
となる．
ここで，$k$は$0 \le k < 1$であるため，$1+k+\cdots+k^{m-n-1} \le 1+k+\cdots+k^{m-1}$となる．
さらに，等比級数から
\begin{equation}
1+k+\cdots+k^{m-1} = \frac{1-k^m}{1-k}
\label{prof:5}
\end{equation}
であるため，(\ref{prof:2})は
\begin{equation}
\| u_n - u_m \| \le \frac{k^n(1-k^m)}{1-k}  \| u_0 - u_1 \|
\label{prof:4}
\end{equation}
となる．
よって，$k$は$0 \le k < 1$から$k^n \rightarrow 0, ~ n \rightarrow \infty$と$k^m \rightarrow 0, ~ m \rightarrow \infty$となる．
即ち，
\begin{align*}
\| u_n - u_m \| \rightarrow 0, ~ (n,m \rightarrow \infty )
\end{align*}
となるため，点列$(u_n)$はCauchy列である．
さらに$X$はBanach空間であるため，$X$は完備である(定義\ref{def:App:Banach})．
即ち，任意のCauchy列が$X$の中で極限を持つ(定義\ref{def:App:Comp})．
よって，点列$(u_n)$は
\begin{align*}
u_n \rightarrow u, ~ n \rightarrow \infty
\end{align*}
となる$u \in X$が存在する．\\
\\
II)

$u_0$を$M$の元とする．
仮定より$A$は$M$から$M$の写像であるため，$A(M) \subset M$となる．
即ち，$u_1 = A u_0$が成立する$u_1 \in M$が存在する．
同様に，$\forall n \in {\mathbb N}$について$u_n \in M$が存在する．
さらに，$M$は閉集合であるため，I)で存在を示した$u$は$M$に属する(定義\ref{def:App:CloseSet})．
よって$Au$も$M$に属する．
その上で仮定より$A$は$k$次の縮小写像であるため，
\begin{align*}
\| A u_n - A u \| \le k \| u_n - u \|
\end{align*}
を得る．
I)より点列$(u_n)$は極限を持つため，$\| u_n - u \| \rightarrow 0, ~ n \rightarrow \infty$となる(定義\ref{def:App:Con})．
即ち
\begin{align*}
\| A u_n - A u \| \rightarrow 0, ~ n \rightarrow \infty
\end{align*}
となるため，$Au$は$Au_n$の極限である(定義\ref{def:App:Con})．
よって，(\ref{pr:Eq:ita})について$n \rightarrow \infty$とすると
\begin{align*}
u=Au
\end{align*}
が成立する．
よって，$u$は(\ref{Eq:FixedPointEq})を満たす真の解$u^*$となる．\\
\\
III)

$u^*,v^* \in M $をそれぞれ$u^* = A u^*$と$v^* = A v^*$を満たすとする．
そのとき，$A$は$k$次の縮小写像であるため，
\begin{align*}
\| u^* - v^* \| = \| A u^* -A v^* \| \le k\|u^* - v^*\|
\end{align*}
を得る．
ここで$0 \le k < 1$であるため，不等式を満たすものは$u^* = v^*$の場合のみである．
即ち，(\ref{Eq:FixedPointEq})を満たす真の解は一意である．

\qed
\end{Proof}

\begin{Proof}[Banachの不動点定理の系]
点列$(u_n)$が収束すること，極限が(\ref{Eq:FixedPointEq})を満たす真の解$u^*$であることは
定理\ref{thm:Banach}の証明のI)とII)と同一の証明で示せている．

(\ref{Ine:Apri})の証明\\
(\ref{prof:4})について$m \rightarrow \infty$とすると$k^m \rightarrow 0$より
\begin{align*}
\| u_n - u \| &\le \frac{k^n}{1-k}  \| u_1 - u_0 \|
\end{align*}

(\ref{Ine:Apos})の証明\\
(\ref{prof:3})より
\begin{align*}
\| u_{n+1} - u_m \| &\le \| u_{n+1} - u_{n+2} \| + \| u_{n+2} - u_{n+3} \| + \cdots + \| u_{m-1} - u_m \|\\
&\le k \| u_n - u_{n+1} \| + k^2 \| u_n - u_{n+1} \| + \cdots + k^{m-n-1} \| u_n - u_{n+1} \| \\
&= k (1 + k + \cdots + k^{m-n-1}) \| u_n - u_{n+1} \|
\end{align*}
を得る．
ここで$1 + k + \cdots + k^{m-n-1} \le 1 + k + \cdots + k^{m-1}$より，
\begin{align*}
\| u_{n+1} - u_m \| &\le k  (1 + k + \cdots + k^{m-1}) \| u_n - u_{n+1} \|
\end{align*}
となる．
よって(\ref{prof:5})より
\begin{align*}
\| u_{n+1} - u_m \| &\le  k \frac{1-k^m}{1-k} \| u_n - u_{n+1} \|
\end{align*}
となる．
よって$m \rightarrow \infty$とすると
\begin{align*}
\| u_{n+1} - u \| &\le  \frac{k}{1-k} \| u_{n+1} - u_n  \|
\end{align*}

(\ref{Ine:rate})の証明\\
\begin{align*}
\| u_{n+1} - u \| = \| Au_n - Au \| \le k \| u_n - u \|
\end{align*}

\qed
\end{Proof}




%----------------------------------------------------------------------------------------------------------------------------------------------
\subsection{方程式の解の品質保証のための基本定理}\label{sec:define_Basic_theorem}
本節では方程式の解の品質保証を行うための基本的な定理を紹介する．
$X, Y$をBanach空間とし，非線形方程式
\begin{align*}
	\mbox{Find} ~ u \in X ~ \mbox{s.t.} ~ F(u) = 0
\end{align*}
を求める問題を考えます．
この問題は有限次元や無限次元を問わず，連立一次方程式や非線形方程式，偏微分方程式などを取り扱うことができます．
この問題に対して近似解の品質保証を行う定理として，以下のNewton-Kantorovichの定理の亜種を紹介します:


\begin{Thm}[基本定理]\label{thm:Basic_Theorem}
	$X$と$Y$をBanach空間とする．
	$F:X \rightarrow Y$を与えられた作用素とし，$\hat{u} \in X$を与えられているとする($\hat{u}$はコンピュータで求めた近似解をイメージで!)．
	$R \in {\mathcal B}(Y, X)$とする．
	$F$は$\hat{u}$でFr\'echet微分可能とし$F'[\hat{u}]$と表記する．
	$F'[\hat{u}]$は全射であるとする．
	$\eta$と$\delta$を不等式
	\begin{align*}
	\| R F(\hat{u}) \|_{X} \le \eta
	\end{align*}
	と
	\begin{align*}
	\| I -  R F'[\hat{u}] \|_{{\mathcal B}(X)} \le \delta < 1
	\end{align*}
	を満たす定数とする．
	
	$\bar{B}(0, 2\eta/(1-\delta))=\{ v \in X ~|~ \| v \|_{X} \le 2\eta/(1-\delta) \}$とする．
	$F$を$\bar{B}(\hat{u}, 2\eta/(1-\delta))$上でFr\'echet微分可能であるとし，$F': \bar{B}(\hat{u}, 2\eta/(1-\delta)) \rightarrow {\mathcal B}(X, Y)$が$\bar{B}(\hat{u}, 2\eta/(1-\delta))$上で連続であるとする．

	$K$を不等式
	\begin{align*}
	\| R \left( F'[\hat{u}] -  F'[\hat{u} + v] \right) \|_{{\mathcal B}(X)} \le K, ~ v \in \bar{B}\left( 0, \frac{2\eta}{1-\delta} \right)
	\end{align*}
	を満たす定数とする.
	もし$2K+\delta \le 1$ならば，
	\begin{align*}
	\| u^* - \hat{u} \|_{X} \le \frac{ \eta }{1 - (K+\delta)} =: \rho
	\end{align*}
	に対し，真の解$u^*$は$\bar{B}(\hat{u}, \rho)$内に存在する.
	その上，$\bar{B}(\hat{u}, 2\eta/(1-\delta))$内で一意である.
\end{Thm}

\begin{Proof}
	まず，$R$と$F'[\hat{u}]$が全単射であることを示す．
	$F'[\hat{u}]$は$F$の$\hat{u}$におけるFr\'echet微分であることから，$F'[\hat{u}]$は${\mathcal B}(X, Y)$に属する．
	さらに，Neumann級数の定理\ref{thm:nuemann_ser}と仮定$\| I -  R F'[\hat{u}] \|_{{\mathcal B}(X)} \le \delta < 1$より，$R F'[\hat{u}]$は全単射である．
	その上，定理\ref{thm:injective_surjective}より，$F'[\hat{u}]$は単射であり，$R$は全射である．
	また，定理の仮定より$F'[\hat{u}]$は全単射となるため，逆作用素$F'[\hat{u}]^{-1}$が存在し，${\mathcal B}(Y, X)$に属する．
	また，$R$の単射性については，定理\ref{thm:linear_operator_injective11}と逆作用素$F'[\hat{u}]^{-1}$を用いて
	\begin{align*}
		\mbox{Find} ~  \phi \in Y ~ \mbox{s.t.} ~ R \phi = 0 \Rightarrow  R F'[\hat{u}] F'[\hat{u}]^{-1} \phi = 0 \Rightarrow \phi = 0
	\end{align*}
	となるため，$R$も全単射である．

	
	続いて，Banachの不動点定理\ref{thm:Banach}を用いて解の存在を示す．
	まず，作用素方程式$F(u) = 0$を不動点方程式に変形する．
	$w := u^* - \hat{u}$とする.
	$R$が単射であるため，
	\begin{align*}
		& F(u^*) = 0 \\
	\Leftrightarrow	& w = w - R F(\hat{u} + w) \\
	\Leftrightarrow	& w = -R F(\hat{u}) +  w - R \left(  F(\hat{u} + w) - F(\hat{u})  \right).
	\end{align*}
	${\mathcal T}: X \rightarrow X$を
	\begin{align*}
		{\mathcal T}(w) := -R F(\hat{u}) +  w - R \left( F(\hat{u} + w) - F(\hat{u}) \right)
	\end{align*}
	となる非線形作用素とし，不動点方程式$w = {\mathcal T}(w)$の解の存在をBanachの不動点定理\ref{thm:Banach}を用いて示す．	
	
	Banachの不動点定理\ref{thm:Banach}では$M$を決めて，${\mathcal T}$が$M$から$M$への縮小写像になることを確認しなければならない．
	特に，ポイントの一つは${\mathcal T}$の定義域を$M$としたときに，値域が$M$に含まれることを確かめなければならない．
	すなわち，${\mathcal T}(M) \subset M$となるように$M$を選ぶことが重要である．
	この定理では，$M = \bar{B}(0, \rho), ~ \rho = \eta /(1 - (K+\delta))$と選ぶ．
	まず，$M$として選んだ閉球$\bar{B}(0, \rho)$と定理で出てくる，もう一つの閉球$\bar{B}(0, 2\eta/(1 - \delta))$の関係性を確認する．	
	定理の仮定より$1 - \delta > 0$と$1 - (\delta + 2 K) \ge 0$を持つため，
	\begin{align*}
		& \rho -  \frac{2\eta}{1 - \delta} = \frac{ \eta (1 - \delta) }{(1 - (\delta + K))(1 - \delta)} - \frac{ 2 \eta (1 - (\delta + K)) }{(1 - (\delta + K))(1 - \delta)} \\
		=& \frac{ \eta \left( (1 - \delta) - 2  (1 - \delta ) + 2 K \right) }{(1 - (\delta +  K))(1 - \delta)} = \frac{ - \eta \left( 1 - \delta - 2 K \right) }{(1 - (\delta + K))(1 - \delta)} \le 0
	\end{align*}
	となる．
	よって，$\rho \le 2\eta/(1 - \delta) $から$\bar{B}(0, \rho) \subset \bar{B}(0, 2\eta/(1 - \delta))$となる．
	
	次に${\mathcal T}( \bar{B}(0, \rho) ) \subset \bar{B}(0, \rho)$を示す.
	仮定から$F'$が$\bar{B}(\hat{u}, 2\eta/(1-\delta))$上で連続であることから，任意の$w \in \bar{B}(0, \rho)$に対しFr\'echet微分とBochner積分に対する微分積分学の基本定理\ref{thm:fundamental_theorem_of_calculus}を用いることで
	\begin{align*}
		\| {\mathcal T}(w) \|_{X} \le& \left\| R F(\hat{u}) \right\|_{X} + \left\| w - R \left(  F(\hat{u} + w) - F(\hat{u})  \right) \right\|_{X} \\
		\le& \eta + \left\| w - R \int_0^1 F'[(1-t) \hat{u} + t (\hat{u} + w) ] w dt \right\|_{X} \\
		\le& \eta + \int_0^1 \left\| w - R F'[\hat{u} + t w ] w \right\|_{X} dt \\
		\le& \eta + \int_0^1 \left\| I - R F'[\hat{u} + t w ]  \right\|_{{\mathcal B}(X)} \| w \|_{X} dt\\
		\le& \eta + \int_0^1 \left( \left\| R (  F'[\hat{u}] -  F'[\hat{u} + t w ])  \right\|_{{\mathcal B}(X)} +  \left\| I - R F'[\hat{u} ]  \right\|_{{\mathcal B}(X)} \right) \| w \|_{X} dt\\
		\le& \eta + \int_0^1 \left( \left\| R (  F'[\hat{u}] -  F'[\hat{u} + t w ])  \right\|_{{\mathcal B}(X)} +  \delta \right) \| w \|_{X} dt
	\end{align*}
	を得る．
	さらに$t \in [0, 1]$に対し$tw \in \bar{B}(0, \rho) \subset \bar{B}(0, 2\eta/(1 - \delta))$となるため，定理の仮定$\left\| R (  F'[\hat{u}] -  F'[\hat{u} + t w ])  \right\|_{{\mathcal B}(X)} \le K$から
	\begin{align*}
		\| {\mathcal T}(w) \|_{X} \le& \eta + \left( K + \delta  \right) \| w \|_{X} \\
		\le& \eta + \left( K + \delta  \right) \rho \\
		=& \eta + \frac{ \eta \left( K + \delta  \right) }{ 1 - ( \delta + K ) } \\
		=& \frac{ \eta - \eta (\delta + K) }{ 1 - ( \delta + K ) } + \frac{ \eta \left( K + \delta  \right) }{ 1 - ( \delta + K ) } = \rho
	\end{align*}
	となる．
	よって，任意の$w \in \bar{B}(0, \rho)$に対して$\| {\mathcal T}(w) \|_{X} \le \rho$となることから，${\mathcal T}( \bar{B}(0, \rho) ) \subset \bar{B}(0, \rho)$となる．
		
	次に${\mathcal T}: \bar{B}(0, \rho) \rightarrow \bar{B}(0, \rho)$が縮小写像になることを確認する．
	すなわち
	\begin{align*}
		\left\| {\mathcal T}(w_1) - {\mathcal T}(w_2) \right\|_{X} \le k \| w_1 - w_2 \|_{X}, ~ \forall w_1, w_2 \in \bar{B}(0, \rho)
	\end{align*}
	となる定数$k$が$1$未満になることを確認しなければならない．
	この定理では，一意性の範囲を広げるために，$\bar{B}(0, \rho) \subset \bar{B}(0, 2\eta/(1 - \delta))$であることを用いて
	\begin{align*}
		\left\| {\mathcal T}(w_1) - {\mathcal T}(w_2) \right\|_{X} \le k \| w_1 - w_2 \|_{X}, ~ \forall w_1, w_2 \in \bar{B}\left(0, \frac{2 \eta}{1 - \delta} \right)
	\end{align*}
	となる定数$k$が$1$未満になることを確かめる．
	その上，任意の$w_1, w_2 \in \bar{B}(0, 2 \eta/(1-\delta))$に対し，
	\begin{align*}
	\left\| {\mathcal T}(w_1) - {\mathcal T}(w_2) \right\|_{X} =& \left\| w_1 - w_2 - R \left( F(\hat{u} + w_1) - F(\hat{u} + w_2) \right) \right\|_{X} \\
		=& \left\| (w_1 - w_2) - R  \int_0^1  F'[\hat{u} + (1-t)w_2 + t w_1] (w_1 - w_2)dt \right\|_{X} \\
		\le& \int_0^1 \left\| I - R F'[\hat{u} + (1-t)w_2 + t w_1] \right\|_{{\mathcal B}(X)} dt \| w_1 - w_2 \|_{X} \\
		\le& \int_0^1 \left( \left\| R \left( F'[\hat{u}] -  F'[\hat{u} + (1-t)w_2 + t w_1] \right) \right\|_{{\mathcal B}(X)} + \delta \right) dt \| w_1 - w_2 \|_{X}
	\end{align*}
	となる．
	任意の$w_1, w_2 \in \bar{B}(0, 2\eta/(1 - \delta) )$と$0 \le t \le 1$に対し，$\| (1-t)w_2 + t w_1 \|_{X} \le (1-t) \| w_2 \|_{X} + t \| w_1 \|_{X} \le 2\eta/(1-\delta)$となるため，
	\begin{align*}
	\left\| {\mathcal T}(w_1) - {\mathcal T}(w_2) \right\|_{X} \le& \left( K + \delta \right) \| w_1 - w_2 \|_{X}
	\end{align*}
	となる．
	その上，仮定$2K + \delta \le 1$かつ$\delta < 1$より$K + \delta < 1$も満たすため，${\mathcal T}$は$\bar{B}(0, \rho)$から$\bar{B}(0, \rho)$への縮小写像となる．
	よって，Banachの不動点定理\ref{thm:Banach}より不動点方程式$w = {\mathcal T}(w)$を満たす解が$\bar{B}(0, \rho)$内に一意に存在する．
	$w = u^* - \hat{u}$であったことを思い出すと，作用素方程式$F(u) = 0$を満たす解が$\bar{B}(\hat{u}, \rho)$内に一意に存在する．


	最後に$\bar{B}(\hat{u}, 2 \eta/(1-\delta))$内で一意になることを示す．
	$\bar{B}(0, 2 \eta/(1-\delta))$内に不動点方程式$w = {\mathcal T}(w)$の解が2つあったとする．
	すなわち2つの解$w_1^*, w_2^* \in \bar{B}(\hat{u}, 2 \eta/(1-\delta))$とし，$w_1^*= {\mathcal T}(w_1^*) $と$w_2^*= {\mathcal T}(w_2^*) $を満たすとする．
	そのとき，
	\begin{align*}
		\left\| w_1^* - w_2^* \right\|_{X} = \left\| {\mathcal T}(w_1^*) - {\mathcal T}(w_2^*) \right\|_{X} \le \left( K + \delta \right) \| w_1^* - w_2^* \|_{X}
	\end{align*}
	となる．
	ここで，$2K + \delta < 1$であるため，不等式を満たすものは$w_1^* = w_2^*$の場合のみである．
	すなわち，不動点方程式$w = {\mathcal T}(w)$を満たす解は$\bar{B}(0, 2 \eta/(1-\delta))$内に一意である．
	よって，作用素方程式$F(u) = 0$を満たす解が$\bar{B}(\hat{u}, 2 \eta/(1-\delta))$内に一意に存在する．
	
\qed
\end{Proof}

定理\ref{thm:Basic_Theorem}では，「$F'[\hat{u}]$は全射」であることと，「$\| I -  R F'[\hat{u}] \|_{{\mathcal B}(X)} \le \delta < 1$」を仮定することで，$F'[\hat{u}]$が全単射であることを証明しています．
有限次元の問題であれば，$R$は$R \approx F'[\hat{u}]^{-1}$とすれば，簡単に作れるので，$\| I -  R F'[\hat{u}] \|_{{\mathcal B}(X)} \le \delta < 1$は，さほど大きな問題にはなりません．
しかしながら，無限次元問題の場合，$R$を作ることすら困難になり始めます．
そのうえで，$\| I -  R F'[\hat{u}] \|_{{\mathcal B}(X)} \le \delta < 1$を満たす保証もありません．
そのために，「$F'[\hat{u}]$は全射」であることと，「$\| I -  R F'[\hat{u}] \|_{{\mathcal B}(X)} \le \delta < 1$」とは別の方法であらかじめ，$F'[\hat{u}]$が全単射であることを証明した上で，定理\ref{thm:Basic_Theorem}を使うことも多々あります．
もちろん，その場合も，定理\ref{thm:Basic_Theorem}の系として次のように簡単に得られます:

\begin{Cor}\label{Cor:Kantorovich}
	$X$と$Y$をBanach空間とする．
	$F:X \rightarrow Y$を与えられた作用素とし，$\hat{u} \in X$を与えられているとする($\hat{u}$はコンピュータで求めた近似解をイメージで!)．
	$F$は$\hat{u}$でFr\'echet微分可能とし$F'[\hat{u}]$と表記する．
	$F'[\hat{u}]$は全単射であるとする．
	$\eta$を不等式
	\begin{align*}
	\| F'[\hat{u}]^{-1} F(\hat{u}) \|_{X} \le \eta
	\end{align*}
	を満たす定数とする．
	
	$\bar{B}(0, 2\eta )=\{ v \in X ~|~ \| v \|_{X} \le 2\eta \}$とする．
	$F$を$\bar{B}(\hat{u}, 2\eta)$上でFr\'echet微分可能であるとし，$F': \bar{B}(\hat{u}, 2\eta) \rightarrow {\mathcal B}(X, Y)$が$\bar{B}(\hat{u}, 2\eta)$上で連続であるとする．
	$K$を不等式
	\begin{align*}
	\| F'[\hat{u}]^{-1} \left( F'[\hat{u}] -  F'[\hat{u} + v] \right) \|_{{\mathcal B}(X)} \le K, ~ v \in \bar{B}\left( 0, 2\eta \right)
	\end{align*}
	を満たす定数とする.
	もし$2K \le 1$ならば，
	\begin{align*}
	\| u^* - \hat{u} \|_{X} \le \frac{ \eta }{1 - K} =: \rho
	\end{align*}
	に対し，真の解$u^*$は$\bar{B}(\hat{u}, \rho)$内に存在する.
	その上，$\bar{B}(\hat{u}, 2\eta)$内で一意である.
\end{Cor}


%
%----------------------------------------------------------------------------------------------------------------------------------------------
%----------------------------------------------------------------------------------------------------------------------------------------------
%----------------------------------------------------------------------------------------------------------------------------------------------
%----------------------------------------------------------------------------------------------------------------------------------------------
